\documentclass[11pt]{article}

\usepackage[a4paper,margin=2.7cm]{geometry}
\usepackage[T1]{fontenc}
\usepackage[utf8]{inputenc}
\usepackage{lmodern}
\usepackage{microtype}
\usepackage{amsmath,amssymb,amsthm,mathtools,bm}
\usepackage{aliascnt}
\usepackage{booktabs,array,longtable}
\usepackage{graphicx}
\usepackage{caption}
\usepackage{subcaption}
\usepackage{enumitem}
\usepackage[backend=biber,style=authoryear,natbib=true,maxcitenames=2,maxbibnames=99]{biblatex}
\addbibresource{references.bib}
\usepackage{xcolor}
\usepackage[hidelinks]{hyperref}
\usepackage[nameinlink,noabbrev]{cleveref}

\setlength{\parindent}{1.2em}
\setlength{\parskip}{0.25em}
\setlist{nosep,leftmargin=2em}

% All environments share the theorem counter so that numbering runs in a single
% sequence.  Plain [theorem] sharing would make cleveref name every cross-
% reference "Theorem", because it keys names off the counter.  aliascnt gives
% each environment a distinct counter name that aliases the shared one, which
% keeps the numbering and lets cleveref recover the correct name.
\newtheorem{theorem}{Theorem}[section]
\newaliascnt{proposition}{theorem}
\newtheorem{proposition}[proposition]{Proposition}
\aliascntresetthe{proposition}
\newaliascnt{corollary}{theorem}
\newtheorem{corollary}[corollary]{Corollary}
\aliascntresetthe{corollary}
\newaliascnt{lemma}{theorem}
\newtheorem{lemma}[lemma]{Lemma}
\aliascntresetthe{lemma}
\newaliascnt{assumption}{theorem}
\newtheorem{assumption}[assumption]{Assumption}
\aliascntresetthe{assumption}
\theoremstyle{definition}
\newaliascnt{definition}{theorem}
\newtheorem{definition}[definition]{Definition}
\aliascntresetthe{definition}
\newaliascnt{remark}{theorem}
\newtheorem{remark}[remark]{Remark}
\aliascntresetthe{remark}
\crefname{proposition}{Proposition}{Propositions}
\crefname{corollary}{Corollary}{Corollaries}
\crefname{lemma}{Lemma}{Lemmas}
\crefname{assumption}{Assumption}{Assumptions}
\crefname{definition}{Definition}{Definitions}
\crefname{remark}{Remark}{Remarks}

\newcommand{\R}{\mathbb{R}}
\newcommand{\E}{\mathbb{E}}
\newcommand{\Pp}{\mathbb{P}}
\newcommand{\Var}{\operatorname{Var}}
\newcommand{\Cov}{\operatorname{Cov}}
\newcommand{\tr}{\operatorname{tr}}
\newcommand{\diag}{\operatorname{diag}}
\newcommand{\op}{o_{\Pp}}
\newcommand{\Op}{O_{\Pp}}
\newcommand{\Normal}{\mathcal{N}}
\newcommand{\ThetaSet}{\boldsymbol{\Theta}}
\newcommand{\thetao}{\theta_0}
\newcommand{\phik}{\varphi_K}
\newcommand{\BS}{Brandt and Santa-Clara}

% ---------------------------------------------------------------------------
% Front-matter metadata.  Do not leave placeholder text in a release build;
% code/make_release.py blocks on it.  \draftlabel is printed only when
% \draftmode is set below.
\newcommand{\authorname}{Diogo Ribeiro}
\newcommand{\affiliation}{ESMAD -- IPP}
\newcommand{\email}{\href{mailto:dfr@esmad.ipp.pt}{\texttt{dfr@esmad.ipp.pt}}}
\newcommand{\orcidid}{0009-0001-2022-7072}
\newcommand{\orcid}{ORCID \href{https://orcid.org/\orcidid}{\texttt{\orcidid}}}

% Uncomment the next line for a draft build; it must stay commented for a release.
%\newcommand{\draftmode}{}

\newcommand{\draftlabel}{%
  \ifdefined\draftmode\\[0.3em]\small Draft for mathematical review\fi}

\newcommand{\papertitle}{Dimensional Asymptotics of Euler-Based Simulated Likelihood for Multidimensional Diffusions}
\newcommand{\paperkeywords}{simulated maximum likelihood; diffusion processes; Euler scheme; transition density; Monte Carlo variance; curse of dimensionality; diffusion bridges; nonregular inference}

\hypersetup{
  pdftitle={\papertitle},
  pdfauthor={\authorname},
  pdfsubject={Asymptotics of Euler-based simulated likelihood for multidimensional diffusions},
  pdfkeywords={\paperkeywords}
}
% ---------------------------------------------------------------------------

\title{\textbf{\papertitle}}
\author{\authorname\\[0.2em]\small\affiliation\\[0.1em]\small\email\quad\textperiodcentered\quad\orcid\draftlabel}
\date{30 July 2026}

\begin{document}
\maketitle

\begin{abstract}
A widely used simulation method for likelihood inference in discretely observed diffusions approximates each transition density by simulating an Euler scheme to one step before the observed endpoint and averaging the final Gaussian density. We show that this endpoint estimator is a shrinking-bandwidth kernel estimator whose effective simulation size is $S/M^{K/2}$, not $S$, where $K$ is the state dimension, $M$ the number of Euler substeps and $S$ the number of simulations. We derive exact Brownian moments, the general local moment law, the corrected density-level normalisation $\sqrt{S/M^{K/2}}$, and the pointwise mean squared error, minimised at $M\asymp S^{2/(K+4)}$. We then give a direct counterexample to the joint transition-density statements of \citet{BrandtSantaClara2002}: along $M=S\to\infty$, a sequence satisfying their rate condition and their own regularity assumptions, the simulated density converges in probability to zero for every $K>2$ although the true density is strictly positive. This refutes the consistency asserted in their Lemma 2 and, since $\sqrt S$ times the centred density then diverges, the central-limit statement of their Lemma 3. The proofs of their parameter-estimator theorems consequently fail, and in four dimensions their rate condition is incompatible with mean-square consistency of the simulated density. We do not prove that the resulting argmax estimator is inconsistent, and we set out what such a proof would require. Issues specific to the exchange-rate application are analysed separately in appendices.
\end{abstract}

\noindent\textbf{Keywords:} simulated maximum likelihood; diffusion processes; Euler scheme; transition density; Monte Carlo variance; curse of dimensionality; diffusion bridges; nonregular inference.

\noindent\textbf{JEL classifications:} C13; C15; C32; G12; G15.

\section{Introduction}

Continuous-time diffusion models are often observed only at discrete dates. Their exact transition densities are rarely available, which makes exact maximum likelihood difficult even when simulation of approximate sample paths is straightforward. One influential response is the simulation method introduced by \citet{Pedersen1995a,Pedersen1995b} and developed and applied by \citet{BrandtSantaClara2002}. The method divides each observation interval into $M$ Euler substeps, simulates the process through the first $M-1$ steps, evaluates the Gaussian density of the final Euler step at the observed endpoint, and averages this density over $S$ simulated paths.

For fixed $M$, the construction is elementary: the sample average converges to the transition density of the $M$-step Euler chain as $S\to\infty$. For fixed $S$ it is also clear that making the final time step smaller turns the Gaussian endpoint density into an increasingly narrow and increasingly high function. The mathematical question is therefore not whether the two separate approximations are individually meaningful. It is whether they remain compatible when $M$ and $S$ increase jointly.

The starting point is not new and we do not claim it. The final Gaussian density is a kernel with bandwidth $h^{1/2}$, where $h=1/M$, so the whole construction is a kernel density estimate built from simulated draws. That reading is due to \citet{MilsteinEtAl2004} and is stated for this estimator by \citet{DetempleGarciaRindisbacher2006}. In $K$ dimensions the kernel has effective volume of order $h^{K/2}$ and height of order $h^{-K/2}$, so only simulations falling in a shrinking neighbourhood of the endpoint contribute materially, and the effective number of such simulations is of order
\[
S h^{K/2}=\frac{S}{M^{K/2}}.
\]
This quantity, rather than $S$ alone, controls the Monte Carlo approximation. That a dimension-dependent rate follows, and that the rate condition of \citet{BrandtSantaClara2002} is inadequate, are also prior results; \Cref{sec:related} sets out exactly what was already known, what this paper adds, and what remains open.

The paper makes six main contributions.

First, we derive exact finite-$M$ moments for the estimator under $K$-dimensional Brownian motion. The Euler approximation is exact in this benchmark, so any failure is caused by the endpoint Monte Carlo construction and not by discretisation bias, boundary behaviour, nonlinear coefficients, or parameter identification.

Second, we construct a direct counterexample to the joint transition-density convergence asserted in Lemma 2 of \citet{BrandtSantaClara2002}. For every $K>2$, choosing $M=S=n$ makes the simulated density at the Brownian origin converge in probability to zero. Standard Brownian motion satisfies the regularity assumptions of that paper exactly, and the sequence satisfies its condition $\sqrt{S}/M\to0$, so the counterexample cannot be excluded by strengthening the hypotheses. The same sequence also refutes the joint central-limit statement of their Lemma 3, since $\sqrt{S_n}(\widehat q_{n,n}-p)\to-\infty$ in probability, so that no finite limit distribution exists. Consequently the published proofs of Theorems 1 and 2 fail, since both rest on those lemmas. In the four-dimensional application, their condition requires $S/M^2\to0$, while the variance calculation requires $S/M^2\to\infty$ for mean-square consistency of the simulated transition density.

We are careful about what this does and does not establish. The counterexample concerns the simulated transition-density estimator at a point. It does not by itself exhibit a sequence along which the simulated-likelihood argmax estimator fails to converge, and we do not claim one. \Cref{sec:scope} states the logical status of each result, and \Cref{sec:argmax_gap} sets out what an argmax counterexample would require.

Third, we establish the generic local moment law for a smooth uniformly elliptic diffusion. If $G_h$ denotes one endpoint-density summand, then for $r>1$,
\[
h^{(r-1)K/2}\E[G_h^r]
\longrightarrow
(2\pi)^{-(r-1)K/2}r^{-K/2}
 p(y\mid x)|V(y)|^{-(r-1)/2}.
\]
In particular, $\Var(G_h)$ is of order $h^{-K/2}$. This gives the corrected triangular-array central-limit normalisation
\[
\sqrt{S h^{K/2}}\,\bigl(\widehat q_{M,S}-q_M\bigr),
\]
not $\sqrt{S}(\widehat q_{M,S}-q_M)$.

Fourth, we combine the $O(M^{-1})$ Euler density bias used by \citet{BrandtSantaClara2002} with the corrected simulation variance. The resulting pointwise mean squared error is
\[
O(M^{-2})+O\!\left(\frac{M^{K/2}}{S}\right),
\]
which is minimised at $M\asymp S^{2/(K+4)}$ and yields the rate $S^{-4/(K+4)}$. The estimator therefore exhibits an explicit curse of dimensionality.

Fifth, we trace the density-level result into the log likelihood. The nonlinearity of the logarithm operates on the effective Monte Carlo size $S/M^{K/2}$. Whenever a standard second-order delta expansion is valid, the leading log-density bias is of order $M^{K/2}/S$, rather than uniformly of order $1/S$. This distinction matters because simulated likelihood sums a log transformation over all observations.

Sixth, we examine the mathematical status of the exchange-rate application. We first give finite-sample scaling diagnostics showing that the ratios corresponding to the paper's own rate conditions are not small at the implemented design, so the implementation does not resemble a point along a sequence in which those conditions operate. The empirical model then contains square-root diffusions that violate the smoothness and nondegeneracy assumptions of the generic theorem, though we show that for the interest rates this violation is numerically inconsequential at the reported estimates. The market-incompleteness state is not: its Feller ratio equals one half for every parameter value, and because the estimator eliminates that state in favour of volatility, every Euler substep requires a recovered square root to remain real. Under the time-varying risk-price specification the variance identity becomes an implicit quadratic in volatility, and at the reported U.S.--U.K. estimates its leading coefficient is negative, so the inverse mapping from the incompleteness state to positive volatility is not globally single-valued without a branch-selection condition. This concerns the inverse map, not the implemented estimator, which takes volatility as primitive and reconstructs incompleteness uniquely. The four-dimensional Brownian correlation matrix is positive definite at the reported estimates and at every model-feasible reference state we examined, and robust to perturbation by the reported standard errors; we find no evidence that the calibrated model reaches an invalid matrix, and the residual point is only that the specification supplies no parameterisation guaranteeing validity globally. Finally, the minimum-incompleteness rule is affine in the selected parameter, so its solution lies on the boundary of the feasible set; we distinguish the constraints the published estimator imposes from those a valid model requires, and together with the binding constraint this yields a hybrid and potentially nonregular estimator not covered by the unconstrained SML theorem.

The aim is not to relitigate the substantive economics of a paper published more than two decades ago. Nor do we infer from a false theorem that every numerical estimate obtained with finite $M$ and $S$ must be poor. The aim is narrower and mathematical: to identify the correct dimensional scale of the endpoint estimator, separate sequential from joint convergence, and determine what the published asymptotic arguments do and do not establish.

\subsection{Relation to existing asymptotic analyses}\label{sec:related}

The estimator analysed here originates with \citet{Pedersen1995a,Pedersen1995b} and is the one implemented by \citet{BrandtSantaClara2002}. A substantial literature has already studied it, and it is important to be exact about what that literature establishes, because several statements that might look like the contribution of this paper are older than it.

\emph{What is already known.} \citet{MilsteinEtAl2004} observe that a simulated transition density of the form $N^{-1}\sum_n\varphi_\delta(\bar X_n-y)$, where $\bar X_n$ are Euler realisations at the terminal time and $\varphi_\delta$ is a Gaussian kernel, is a special case of a Parzen--Rosenblatt estimator with bandwidth $\delta$ set equal to the integration step; they then construct forward--reverse representations that are root-$N$ consistent and so avoid the resulting curse of dimensionality. Two features of their setting differ from the one studied here, and the difference is where the exponent below comes from. Their kernel is an \emph{imposed} smoother of bandwidth $\delta=h$, hence of variance $h^2$, applied to paths simulated all the way to the terminal time; the endpoint estimator of \citet{Pedersen1995a,Pedersen1995b} instead stops one step short and uses the model's own one-step Gaussian density, whose covariance is $hV$ and whose bandwidth is therefore $\sqrt h$. The effective sample sizes are correspondingly $N\delta^K=Nh^K$ in their case and $Sh^{K/2}$ here. The dimension-dependent rate they quote, an optimal bandwidth of order $N^{-1/(4+K)}$ and accuracy $N^{-2/(4+K)}$, is the classical Parzen--Rosenblatt rate, which they attribute to the nonparametric-statistics literature.

\citet{DetempleGarciaRindisbacher2006} restate the kernel interpretation for the estimator of \citet{Pedersen1995a,Pedersen1995b} and \citet{BrandtSantaClara2002} specifically, and draw three consequences from it. First, the convergence rate is dimension dependent: they give $M^{-2/(d+4)}$ for the score, where $d$ is the state dimension, which is the same exponent family as the pointwise rate derived below. Second, optimality requires a joint design linking sample size, simulation size and discretisation level, and they state the conditions explicitly. Third, and most directly, they observe that \citet{BrandtSantaClara2002} impose one of those two conditions and not the other, and write that this ``assumption, unfortunately, is not sufficient''. They also note that with a fixed discretisation level and a fixed number of simulations the second-order bias explodes as the sample grows, so that asymptotic confidence intervals cover the true value with probability zero.

Separately, \citet{DurhamGallant2002} observe that unconditioned Euler paths rarely arrive near the observed endpoint and introduce bridge-based proposals for that reason, and \citet{ElerianChibShephard2001} augment the path with a Metropolis--Hastings block on the same reasoning. \citet{Lindstrom2012} regularises the Pedersen proposal towards a bridge, and \citet{StramerYan2007b} compare the endpoint estimator numerically with a closed-form expansion.

\citet{StramerYan2007a} require particular care, because although their theorem concerns a different sampler their remarks about this one are pointed. Their Theorem 1 gives an efficient allocation for the \emph{modified Brownian bridge} sampler, showing that its importance-weight variance is bounded uniformly in $M$ and that the optimal choice is $N\asymp M^2$ with total error $O(1/M)+O(1/\sqrt N)$; they note that this accuracy holds ``regardless of the dimension $d$''. About the endpoint sampler they observe three things. It ``may result a random variable $R_M$ with unbounded $\Var_q(R_M)$ which leads to inefficiency''; in its case ``a limiting quantity for $G_M$ does not exist and, therefore, even a huge number $N$ can give poor approximations''; and numerically ``for any fixed $N$, the approximation error of Pedersen's sampler is dominated by the Monte Carlo error and increases as $M$ increases''. They also record the kernel connection for this estimator specifically, identifying it with a kernel method in which the evaluation points are the preterminal states at time $\Delta-h$.

We therefore claim \emph{none} of the following as new: the kernel interpretation of the endpoint density; the dimension dependence of the resulting rate, including the exponent $2/(d+4)$; the existence of a bias--variance trade-off between discretisation and simulation and the associated optimal allocation; the observation that the rate condition of \citet{BrandtSantaClara2002} is inadequate; the need for joint rather than sequential limits; the case for bridge proposals; the possibility that the endpoint estimator's importance-weight variance is unbounded in $M$; and the observation that increasing $S$ alone need not rescue it.

One point of detail matters for what follows, because the dimensional exponent turns on it. The kernel in the endpoint construction is the model's own one-step Euler density, whose covariance is $hV$; its width is therefore $\sqrt h$, not $h$. That is why the effective simulation size is $Sh^{K/2}$ rather than $Sh^K$, and it is the reason a dimension-dependent exponent appears here where \citet{StramerYan2007a} find none for the bridge.

\emph{What this paper adds.} Against that background the contributions are the following, each stated for the unconditioned endpoint estimator.
\begin{enumerate}[label=(\alph*)]
\item Exact finite-$M$ moments of every order for the endpoint summand under $K$-dimensional Brownian motion (\Cref{prop:brownian_moments}). The prior literature is asymptotic; these are closed forms valid at every $M$ and $S$, and they are what make the counterexample below computable rather than merely plausible.
\item The exact local moment constant $c_{r,K}=(2\pi)^{-(r-1)K/2}r^{-K/2}$ for general uniformly elliptic diffusions (\Cref{thm:local_moments}), together with the identification of $R_{M,S}=S/M^{K/2}$ as the governing quantity. Prior work gives the exponent; this gives the constant.
\item An explicit collapse-in-probability sequence (\Cref{thm:collapse}). The prior negative results are of a different kind: inefficiency, exploding second-order bias in the estimating function, possibly unbounded importance-weight variance, and the warning that a huge $N$ may not help. Each says the estimator may be bad. \Cref{thm:collapse} exhibits a specific sequence, admissible under the published rate condition and the published regularity assumptions, along which the density estimator converges in probability to a value it should not, with the true density strictly positive. That the variance may be unbounded is prior work; that there is an admissible path to the wrong limit is not.
\item A direct refutation of the joint transition-density statements in Lemmas 2 and 3 of \citet{BrandtSantaClara2002} (\Cref{thm:collapse} and \Cref{cor:clt_failure}), under those authors' own assumptions.
\item The explicit four-dimensional incompatibility between the published condition, equivalent to $S/M^2\to0$, and mean-square consistency, which requires $S/M^2\to\infty$. \citet{DetempleGarciaRindisbacher2006} state that the published condition is insufficient; the contribution here is to show that in the dimension actually used no admissible sequence exists at all.
\item The mathematical reassessment of the exchange-rate application in \Cref{sec:application} and \Cref{app:sqrt} to \Cref{app:constrained}.
\end{enumerate}

The right description of this paper's relation to \citet{DetempleGarciaRindisbacher2006} is therefore sharpening rather than discovery: it derives exact finite-discretisation moments, identifies the local moment constants, and constructs an explicit admissible sequence along which the simulated transition density converges to the wrong value. What distinguishes that last item from an efficiency observation is its logical force. An inefficiency says a design is wasteful and can be repaired by spending more. A collapse along an admissible sequence says the density estimator converges to a value it should not; divergence of $S$ alone does not repair it, and only a fast enough rate of $S$ relative to $M^{K/2}$ does.

\Cref{tab:related} summarises the comparison.

\begin{table}[t]
\centering
\caption{Relation to existing asymptotic analyses of Pedersen-type simulated likelihood. ``Joint'' asks whether limits are taken jointly rather than sequentially; ``Counterexample'' asks whether a sequence is exhibited along which the estimator converges to a wrong value.}
\label{tab:related}
\footnotesize
\setlength{\tabcolsep}{4pt}
\newcolumntype{L}[1]{>{\raggedright\arraybackslash}p{#1}}
\begin{tabular}{L{2.1cm}L{2.3cm}L{1.5cm}L{2.9cm}L{1.3cm}L{1.3cm}}
\toprule
Reference & Estimator & Dim. & Rate or variance result & Joint & Counter\-example \\
\midrule
\citet{Pedersen1995a,Pedersen1995b}
& Endpoint & general & consistency and asymptotic normality & seq. & no \\
\citet{BrandtSantaClara2002}
& Endpoint & applied at $K=4$ & $\sqrt S/M\to0$, $N/S^{1/4}\to0$ & asserted, proved seq. & no \\
\citet{MilsteinEtAl2004}
& Terminal-time kernel; forward--reverse & explicit in $K$ & kernel reading with bandwidth $\delta=h$; classical Parzen--Rosenblatt rate; root-$N$ forward--reverse alternative & yes & no \\
\citet{DetempleGarciaRindisbacher2006}
& Euler and Doss-trans\-formed, incl.\ Pedersen and Brandt--Santa-Clara & explicit in $d$ & score rate $M^{-2/(d+4)}$; optimal design; published condition ``not sufficient''; second-order bias correction & yes & no \\
\citet{DurhamGallant2002}
& Bridge proposals & general & numerical efficiency comparisons & --- & no \\
\citet{StramerYan2007a}
& Modified Brownian bridge; remarks on the endpoint sampler & dimension-free for the bridge & bridge weight variance bounded in $M$, optimal $N\asymp M^2$; endpoint variance may be unbounded and large $N$ may not rescue it & yes & no \\
\citet{Lindstrom2012}
& Regularised bridge & general & sparse-sampling performance & --- & no \\
\midrule
This paper
& Endpoint & explicit in $K$ & exact finite-$M$ moments; exact constant $c_{r,K}$; $\Var\asymp M^{K/2}/S$ & yes & yes, \Cref{thm:collapse} \\
\bottomrule
\end{tabular}
\end{table}

\subsection{Scope of the results}\label{sec:scope}

Because this paper both proves theorems and reassesses a published application, it matters which claims are which. Every result below falls into one of three categories, and the text signals the category wherever a result is stated.

\emph{Proved.} The exact Brownian moments of \Cref{prop:brownian_moments}, the collapse theorem \Cref{thm:collapse}, the general local moment law \Cref{thm:local_moments}, the density variance scale of \Cref{cor:generic_variance}, the density-level triangular-array central limit theorems \Cref{thm:brownian_clt} and \Cref{thm:generic_clt}, and the pointwise mean-squared-error balance of \Cref{prop:optimal_m}. These are theorems about the simulated \emph{transition density}, proved under stated assumptions.

\emph{Conditional.} The log-density expansion \Cref{prop:log_delta}, which requires uniform integrability and lower-tail control that we assume rather than verify; the resulting statements about scores and Hessians; and the abstract perturbation result \Cref{prop:m_estimator}, which is a deterministic implication and not a theorem about the simulated estimator.

\emph{Diagnostic.} Everything in \Cref{sec:application} concerning the exchange-rate model: the finite-sample scaling ratios, the square-root boundary behaviour, correlation-matrix feasibility, the volatility state transformation, minimum-incompleteness geometry, and constrained inference. These are calculations about a specific published model, not asymptotic theorems, and they are reported as such.

One boundary deserves emphasis because it is easy to blur. All the proved results concern the simulated transition density. None of them concerns the maximiser of the simulated log likelihood. We do not prove that the endpoint simulated-likelihood estimator is inconsistent, and \Cref{sec:argmax_gap} and \Cref{app:argmax} set out precisely what such a proof would require.

The remainder is organised as follows. \Cref{sec:estimator} defines the estimator. \Cref{sec:brownian} gives exact Brownian results and the counterexample. \Cref{sec:general} derives the general diffusion moment law and corrected CLT. \Cref{sec:mse} studies the bias--variance trade-off. \Cref{sec:loglik} discusses log likelihood and parameter estimation. \Cref{sec:proof_review} analyses the published proof. \Cref{sec:application} studies the exchange-rate diffusion. \Cref{sec:numerics} presents reproducible numerical evidence, and \Cref{sec:remedies} discusses variance-reduction alternatives.

\section{The endpoint simulated-likelihood estimator}\label{sec:estimator}

Let $Y_t\in\R^K$ satisfy the stochastic differential equation
\begin{equation}\label{eq:sde}
 dY_t=\mu(Y_t,t;\theta)\,dt+\Sigma(Y_t,t;\theta)\,dW_t,
 \qquad V=\Sigma\Sigma^{\mathsf T},
\end{equation}
where $W_t$ is a Brownian motion of suitable dimension and $\theta\in\ThetaSet\subset\R^L$. Coefficients are assumed regular enough for a unique strong solution in the sense of \citet{KaratzasShreve1991}. We suppress $t$ and $\theta$ where no ambiguity arises. Consider one observation interval of length one, from $Y_0=x$ to $Y_1=y$. The normalisation to unit length is only notational.

For $M\ge2$, set $h=1/M$ and define the Euler chain \citep{KloedenPlaten1992}
\begin{equation}\label{eq:euler}
Y^M_{m+1}=Y^M_m+\mu(Y^M_m,mh)h
 +\Sigma(Y^M_m,mh)\sqrt{h}\,\varepsilon_{m+1},
\qquad m=0,\ldots,M-1,
\end{equation}
with $Y^M_0=x$ and independent $\varepsilon_m\sim\Normal(0,I)$. Conditional on $Y^M_m=z$, the one-step density is
\begin{equation}\label{eq:one_step_kernel}
g_h(y\mid z)
=
\phik\!\left(y;
 z+\mu(z,1-h)h,
 hV(z,1-h)
\right).
\end{equation}
Let $Z_h=Y^M_{M-1}$ denote the preterminal Euler state and let $r_h(z\mid x)$ be its density. The $M$-step Euler transition density is
\begin{equation}\label{eq:euler_density_expectation}
q_M(y\mid x)=\E[g_h(y\mid Z_h)]
=\int_{\R^K}g_h(y\mid z)r_h(z\mid x)\,dz.
\end{equation}

The endpoint simulation estimator is
\begin{equation}\label{eq:pedersen_estimator}
\widehat q_{M,S}(y\mid x)
=
\frac{1}{S}\sum_{s=1}^S g_h(y\mid Z_{h,s}),
\end{equation}
where $Z_{h,1},\ldots,Z_{h,S}$ are independent draws from the preterminal Euler law. This is equation (8) of \citet{BrandtSantaClara2002}, with notation adapted to the present analysis.

For each fixed $M$, the strong law yields
\[
\widehat q_{M,S}(y\mid x)\longrightarrow q_M(y\mid x)
\quad\text{almost surely as }S\to\infty,
\]
provided the summand is integrable. Under regularity conditions such as those of \citet{BallyTalay1996a,BallyTalay1996b},
\[
q_M(y\mid x)-p(y\mid x)=O(M^{-1}),
\]
where $p$ is the diffusion transition density. These two facts establish a sequential limit: first $S\to\infty$ for fixed $M$, then $M\to\infty$. They do not establish convergence along arbitrary joint sequences $M,S\to\infty$.

The geometry of \eqref{eq:one_step_kernel} explains why. Its covariance is $hV(z)$. Under nondegeneracy, it is concentrated in a ball of radius $O(\sqrt h)$ around the endpoint and has maximum height $O(h^{-K/2})$. The integral in \eqref{eq:euler_density_expectation} remains finite because the effective volume is $O(h^{K/2})$. A Monte Carlo average resolves that integral only if enough preterminal draws enter this shrinking region.

\begin{definition}[Effective endpoint simulation size]\label{def:effective_size}
For a $K$-dimensional endpoint estimator with Euler step $h=1/M$, define
\begin{equation}\label{eq:effective_size}
R_{M,S}=S h^{K/2}=\frac{S}{M^{K/2}}.
\end{equation}
This is the expected-order number of simulations falling in an $O(\sqrt h)$ neighbourhood of an endpoint with a regular positive preterminal density.
\end{definition}

The next sections make this interpretation exact.

\section{Exact Brownian analysis}\label{sec:brownian}

Consider standard Brownian motion in $\R^K$,
\begin{equation}\label{eq:brownian}
dY_t=dW_t,
\end{equation}
observed over a unit interval. Euler discretisation is exact. This benchmark is not an adversarial construction outside the scope of the published theory: with $\mu\equiv0$ and $\Sigma\equiv I_K$, all coefficient derivatives vanish and are therefore continuous and bounded of every order, so Assumption 1 of \citet{BrandtSantaClara2002} holds, and $\Sigma\Sigma^{\mathsf T}=I_K$ is positive definite, so Assumption 2 holds. Lemmas 1--3 of that paper are stated under exactly these two assumptions. Any failure exhibited here is therefore a failure inside the stated hypotheses, and cannot be repaired by strengthening the regularity conditions. Set $Y_0=x$ and evaluate the transition density at $Y_1=y$. With $h=1/M$,
\[
Z_h\sim\Normal(x,(1-h)I_K)
\]
and a single endpoint summand is
\begin{equation}\label{eq:brownian_summand}
G_h
=(2\pi h)^{-K/2}
\exp\!\left[-\frac{\|y-Z_h\|^2}{2h}\right].
\end{equation}

\subsection{Exact moments}

\begin{proposition}[Exact Brownian moments]\label{prop:brownian_moments}
Let $r>0$. For the summand \eqref{eq:brownian_summand},
\begin{align}
\E[G_h^r]
&=(2\pi h)^{-rK/2}
\left(1+\frac{r(1-h)}{h}\right)^{-K/2}
\exp\!\left[
-\frac{r\|y-x\|^2}{2\{h+r(1-h)\}}
\right] \label{eq:brownian_r_moment}\\
&=(2\pi)^{-rK/2}
 h^{-(r-1)K/2}
\{r-(r-1)h\}^{-K/2}
\exp\!\left[
-\frac{r\|y-x\|^2}{2\{r-(r-1)h\}}
\right]. \nonumber
\end{align}
In particular,
\begin{equation}\label{eq:brownian_mean}
\E[G_h]=(2\pi)^{-K/2}\exp\!\left[-\frac{\|y-x\|^2}{2}\right]
=p(y\mid x),
\end{equation}
and
\begin{equation}\label{eq:brownian_second}
\E[G_h^2]
=(2\pi)^{-K}h^{-K/2}(2-h)^{-K/2}
\exp\!\left[-\frac{\|y-x\|^2}{2-h}\right].
\end{equation}
\end{proposition}

\begin{proof}
For $Z\sim\Normal(x,(1-h)I_K)$ and $c>0$, the Gaussian identity
\[
\E\!\left[e^{-c\|Z-y\|^2}\right]
=\{1+2c(1-h)\}^{-K/2}
\exp\!\left[-\frac{c\|x-y\|^2}{1+2c(1-h)}\right]
\]
follows by completing the square. Apply it with $c=r/(2h)$ and multiply by the prefactor $(2\pi h)^{-rK/2}$. The cases $r=1$ and $r=2$ follow directly.
\end{proof}

\begin{corollary}[Exact variance]\label{cor:brownian_variance}
For $S$ independent simulations,
\begin{equation}\label{eq:brownian_estimator_variance}
\Var(\widehat q_{M,S})
=
\frac{1}{S}\left[
(2\pi)^{-K}h^{-K/2}(2-h)^{-K/2}
 e^{-\|y-x\|^2/(2-h)}
-p(y\mid x)^2
\right].
\end{equation}
Hence
\begin{equation}\label{eq:brownian_variance_asymptotic}
\Var(\widehat q_{M,S})
\sim
\frac{\tau_K^2(x,y)}{S h^{K/2}},
\qquad
\tau_K^2(x,y)
=(2\pi)^{-K}2^{-K/2}e^{-\|y-x\|^2/2}.
\end{equation}
\end{corollary}

The variance is not uniformly $O(1/S)$. It is $O(M^{K/2}/S)$. The estimator is unbiased in this example for every $M$, but its distribution becomes increasingly dominated by rare, very large contributions when $R_{M,S}$ is small.

\subsection{A direct joint-limit counterexample}

A variance calculation alone does not prove failure of convergence in probability because a non-uniformly integrable sequence can converge in probability while retaining large or divergent moments. The next theorem gives a direct failure result. It is a statement about the simulated transition density at a single pair of points; \Cref{sec:argmax_gap} explains why it does not transfer automatically to the parameter estimator.

\begin{theorem}[Collapse along $M=S$]\label{thm:collapse}
Let $K>2$, $x=y=0$, and choose $M_n=S_n=n$. Then the Brownian endpoint estimator satisfies
\begin{equation}\label{eq:collapse}
\widehat q_{n,n}(0\mid0)\longrightarrow0
\qquad\text{in probability},
\end{equation}
although
\[
p(0\mid0)=(2\pi)^{-K/2}>0.
\]
In every dimension the sequence satisfies the rate condition
\[
\frac{\sqrt{S_n}}{M_n}=n^{-1/2}\longrightarrow0
\]
used in Theorems 1 and 2 of \citet{BrandtSantaClara2002}, and it satisfies without any rate restriction the hypotheses of their Lemma 2, which requires only $M\to\infty$ and $S\to\infty$.
\end{theorem}

\begin{proof}
Write $h_n=1/n$ and let $Z_{n,1},\ldots,Z_{n,n}$ be the preterminal draws, each distributed as $\Normal(0,(1-1/n)I_K)$. Fix a constant $c>K$ and define
\[
r_n=\sqrt{\frac{c\log n}{n}}.
\]
Let $A_n$ be the event that at least one draw lies in the ball $B(0,r_n)$. The Gaussian densities of the $Z_{n,s}$ are uniformly bounded for $n\ge2$, so for a constant $C_K$,
\[
\Pp(A_n)
\le n\Pp(\|Z_{n,1}\|\le r_n)
\le C_K n r_n^K
=C_Kc^{K/2}n^{1-K/2}(\log n)^{K/2}
\longrightarrow0
\]
because $K>2$.

On $A_n^c$, every summand is bounded by
\[
\left(\frac{n}{2\pi}\right)^{K/2}
\exp\!\left(-\frac{nr_n^2}{2}\right)
=(2\pi)^{-K/2}n^{(K-c)/2},
\]
which tends to zero because $c>K$. The sample average obeys the same bound. Therefore, for every $\varepsilon>0$,
\[
\Pp(\widehat q_{n,n}>\varepsilon)
\le \Pp(A_n)+
\mathbf{1}\!\left\{(2\pi)^{-K/2}n^{(K-c)/2}>\varepsilon\right\}
\longrightarrow0.
\]
\end{proof}

The same sequence refutes a second published statement. Lemma 3 of \citet{BrandtSantaClara2002} asserts that, under $\sqrt S/M\to0$,
\[
\sqrt S\left\{\widehat q_{M,S}(y\mid x)-p(y\mid x)\right\}
\]
converges in distribution to a normal law. That cannot hold along $M_n=S_n=n$.

\begin{corollary}[Failure of the published joint central-limit statement]\label{cor:clt_failure}
Let $K>2$ and $x=y=0$, and take $M_n=S_n=n$, a sequence satisfying $\sqrt{S_n}/M_n\to0$. Then
\begin{equation}\label{eq:clt_divergence}
\sqrt{S_n}\left\{\widehat q_{n,n}(0\mid0)-p(0\mid0)\right\}\longrightarrow-\infty
\qquad\text{in probability},
\end{equation}
so the sequence converges in distribution to no finite limit, normal or otherwise. The joint central-limit statement of Lemma 3 is therefore false along an admissible sequence, whatever variance is assigned to the limit.
\end{corollary}

\begin{proof}
Write $p=p(0\mid0)=(2\pi)^{-K/2}>0$ and fix $\Lambda>0$. By \Cref{thm:collapse}, $\widehat q_{n,n}\to0$ in probability, so with $\varepsilon=p/2$,
\[
\Pp\left(\widehat q_{n,n}>\tfrac p2\right)\longrightarrow0 .
\]
On the complementary event, $\widehat q_{n,n}-p\le-\tfrac p2$, hence $\sqrt{S_n}(\widehat q_{n,n}-p)\le-\tfrac p2\sqrt n$. Since $\tfrac p2\sqrt n>\Lambda$ for all $n$ large,
\[
\Pp\left(\sqrt{S_n}\left\{\widehat q_{n,n}-p\right\}>-\Lambda\right)
\le\Pp\left(\widehat q_{n,n}>\tfrac p2\right)\longrightarrow0 ,
\]
which is \eqref{eq:clt_divergence}. A sequence diverging to $-\infty$ in probability is not tight, so it has no weak limit in $\R$.
\end{proof}

\begin{remark}[What is and is not refuted]\label{rem:hierarchy}
It is worth stating the hierarchy in one place, since the four statements are easy to run together.
\begin{enumerate}[label=(\roman*)]
\item The joint transition-density consistency of Lemma 2 is \emph{directly false} (\Cref{thm:collapse}).
\item The joint central-limit statement of Lemma 3 is \emph{directly false} along the same sequence (\Cref{cor:clt_failure}). Note also that the limiting variance it states is the variance of a summand that itself depends on $M$, so the statement is ill-posed as an asymptotic claim even before the counterexample is applied.
\item The published proofs of Theorems 1 and 2 \emph{fail}, since Lemmas 4 to 6 take Lemmas 2 and 3 as inputs.
\item The conclusions of the parameter-estimator theorems are \emph{not} disproved. The construction says nothing directly about the simulated-likelihood argmax; see \Cref{sec:argmax_gap} and \Cref{app:argmax}.
\end{enumerate}
\end{remark}

\begin{remark}[How an unbiased estimator converges to zero]
For every $n$, $\E[\widehat q_{n,n}]=p(0\mid0)$. The convergence in probability to zero is sustained by increasingly rare simulations with increasingly large endpoint densities. The sequence is not uniformly integrable. This is precisely the rare-event geometry hidden by a fixed-$M$ application of the strong law.
\end{remark}

\begin{corollary}[Contradiction in four dimensions]\label{cor:contradiction}
For $K=4$, mean-square consistency requires
\[
\frac{S}{M^2}\longrightarrow\infty,
\]
whereas the published condition $\sqrt S/M\to0$ is equivalent to
\[
\frac{S}{M^2}\longrightarrow0.
\]
The two conditions are mutually exclusive. Moreover, \Cref{thm:collapse} shows that the published condition permits a sequence along which the density estimator converges to the wrong value.
\end{corollary}

\begin{remark}[Where the conflict is unconditional]
The incompatibility in \Cref{cor:contradiction} is a property of dimensions $K\ge4$. Mean-square consistency requires $S/M^{K/2}\to\infty$ and the published condition requires $S/M^2\to0$; for $K\ge4$ the first demands $S\gtrsim M^{K/2}\ge M^2$ and the second demands $S\ll M^2$, so no sequence satisfies both. For $K\le3$ the two are compatible: with $K=3$ and $S=M^{7/4}$ one has $S/M^{3/2}\to\infty$ and $\sqrt S/M\to0$ simultaneously. \Cref{thm:collapse} is nevertheless not confined to $K\ge4$. It shows that even when the published condition admits consistent sequences, as at $K=3$, it also admits inconsistent ones, so the condition cannot be what makes the theorem work. The empirical application of \citet{BrandtSantaClara2002} has $K=4$, where no admissible sequence exists at all.
\end{remark}

\subsection{Correct triangular-array central limit theorem}

\begin{theorem}[Brownian endpoint CLT]\label{thm:brownian_clt}
Let $M\to\infty$ and $S\to\infty$ such that
\begin{equation}\label{eq:effective_diverge}
S h^{K/2}=\frac{S}{M^{K/2}}\longrightarrow\infty.
\end{equation}
Then
\begin{equation}\label{eq:brownian_clt}
\sqrt{S h^{K/2}}
\left\{\widehat q_{M,S}(y\mid x)-p(y\mid x)\right\}
\Longrightarrow
\Normal\!\left(0,\tau_K^2(x,y)\right),
\end{equation}
where $\tau_K^2(x,y)$ is defined in \eqref{eq:brownian_variance_asymptotic}.
\end{theorem}

\begin{proof}
Index the array by $n$, with $h_n=1/M_n$ and $s=1,\ldots,S_n$; within row $n$ the summands $G_{n,s}$ are independent and identically distributed. \Cref{prop:brownian_moments} at $r=2$ gives $h_n^{K/2}\E[G_{n,1}^2]\to\tau_K^2$, and $\E G_{n,1}=p(y\mid x)$ exactly for every $M$, so
\[
h_n^{K/2}\Var(G_{n,1})=h_n^{K/2}\E[G_{n,1}^2]-h_n^{K/2}p(y\mid x)^2\longrightarrow\tau_K^2 .
\]
At $r=3$ the same proposition gives $\E[G_{n,1}^3]=O(h_n^{-K})$, a bound on the \emph{raw} moment. Since $G_{n,1}\ge0$ and $\E G_{n,1}=p(y\mid x)$ is constant in $n$, the inequality $|a-b|^3\le4(a^3+b^3)$ for $a,b\ge0$ gives the centred bound
\[
\E\lvert G_{n,1}-\E G_{n,1}\rvert^3\le4\left\{\E[G_{n,1}^3]+p(y\mid x)^3\right\}=O(h_n^{-K}).
\]
The Lyapunov ratio at order three is therefore
\[
\frac{S_n\E\lvert G_{n,1}-\E G_{n,1}\rvert^3}
{\{S_n\Var(G_{n,1})\}^{3/2}}
=\frac{S_nO(h_n^{-K})}{S_n^{3/2}\tau_K^3h_n^{-3K/4}\{1+o(1)\}}
=O\!\left(\left(S_nh_n^{K/2}\right)^{-1/2}\right)\longrightarrow0
\]
by \eqref{eq:effective_diverge}. The triangular-array Lyapunov central limit theorem applies. Because the Euler scheme is exact for Brownian motion, $q_{M}=p$ for every $M$ and no bias term appears, so the statement may be centred at $p$ directly.
\end{proof}

The correct normalisation depends on dimension. A fixed-$M$ $\sqrt S$ central limit theorem is valid for each $M$ separately, but it is not stable as $M\to\infty$: see \Cref{rem:not_uniform}.

\section{General multidimensional diffusions}\label{sec:general}

The Brownian calculation reflects a local property of Gaussian endpoint kernels. We now derive the generic moment law. We state it under a clean global framework, so that every step of the proof is checkable and no condition is phrased in terms of the conclusion it is meant to deliver; \Cref{rem:localisation} explains how localisation would weaken it.

Throughout this section the model is time homogeneous, so that $\mu(z)$ and $V(z)$ carry no time argument. The one-step kernel \eqref{eq:one_step_kernel} is then evaluated at the preterminal state alone, and no ambiguity arises about whether coefficients are taken at time $1-h$ or at time $1$. \Cref{rem:time_inhomogeneous} records what changes otherwise.

Let $G_h=g_h(y\mid Z_h)$ be one summand in \eqref{eq:pedersen_estimator}, where $Z_h$ is the Euler chain at time $1-h$. Fix $x$ and $y$.

\begin{assumption}[Global regularity]\label{ass:local}
\begin{enumerate}[label=(A\arabic*)]
\item\label{ass:mu} $\mu:\R^K\to\R^K$ is bounded and continuous, with $\lVert\mu\rVert_\infty=\bar\mu<\infty$.
\item\label{ass:V} $V:\R^K\to\R^{K\times K}$ is continuous and globally uniformly elliptic: there exist constants $0<\underline v\le\overline v<\infty$ with
\begin{equation}\label{eq:uniform_ellipticity}
\underline v\,I_K\preceq V(z)\preceq\overline v\,I_K
\qquad\text{for every }z\in\R^K.
\end{equation}
\item\label{ass:rh} For each $h\in(0,\tfrac12]$ the preterminal Euler state $Z_h$, that is the chain \eqref{eq:euler} at time $1-h$, has a density $r_h(\cdot\mid x)$ with respect to Lebesgue measure.
\item\label{ass:bound} The family is uniformly bounded: $\sup_{0<h\le1/2}\ \sup_{z\in\R^K}r_h(z\mid x)=\bar r<\infty$.
\item\label{ass:conv} There is a continuous density $p(\cdot\mid x)$ with $p(y\mid x)>0$ such that $r_h(\cdot\mid x)\to p(\cdot\mid x)$ uniformly on some closed ball $\overline B(y,\delta_0)$, $\delta_0>0$, as $h\downarrow0$.
\end{enumerate}
\end{assumption}

Two comments on the content of these conditions. \ref{ass:conv} is a joint requirement in time and space: the preterminal state lives at time $1-h$, not at time $1$, so the stated convergence combines convergence of the Euler law to the diffusion law with continuity of $t\mapsto p_t(\cdot\mid x)$ at $t=1$ in the topology of uniform convergence near $y$. The limit $p(\cdot\mid x)$ is the diffusion transition density at time $1$. And \ref{ass:bound} is not implied by \ref{ass:rh}: it is a uniform-in-$h$ statement, which for uniformly elliptic coefficients follows from the Gaussian upper bounds for Euler densities established by \citet{BallyTalay1996b}, since the preterminal time $1-h\ge\tfrac12$ is bounded away from zero.

\Cref{ass:local} is compatible with the regular setting used to justify Euler density expansions in \citet{BallyTalay1996a,BallyTalay1996b}. It is not compatible with the square-root boundaries in the later financial application, where \ref{ass:V} fails at the origin; that distinction is discussed in \Cref{sec:application}.

The following theorem proves the dimensional moment law under \Cref{ass:local}. It concerns the simulated transition density.

\begin{theorem}[Local moment law]\label{thm:local_moments}
Under \Cref{ass:local}, for every fixed $r>1$,
\begin{equation}\label{eq:general_r_moment}
 h^{(r-1)K/2}\E[G_h^r]
\longrightarrow
c_{r,K}\,p(y\mid x)\,|V(y)|^{-(r-1)/2},
\end{equation}
where
\begin{equation}\label{eq:crk}
c_{r,K}=(2\pi)^{-(r-1)K/2}r^{-K/2}.
\end{equation}
In particular,
\begin{equation}\label{eq:general_second}
 h^{K/2}\E[G_h^2]
\longrightarrow
(4\pi)^{-K/2}p(y\mid x)|V(y)|^{-1/2}.
\end{equation}
\end{theorem}

\begin{proof}
Write $m_h(z)=y-z-h\mu(z)$, so that by \eqref{eq:one_step_kernel}
\begin{equation}\label{eq:Ghr}
g_h(y\mid z)^r
=(2\pi h)^{-rK/2}|V(z)|^{-r/2}
\exp\!\left[-\frac{r}{2h}m_h(z)^{\mathsf T}V(z)^{-1}m_h(z)\right],
\end{equation}
and fix $\delta\in(0,\delta_0]$. Split
\begin{align*}
h^{(r-1)K/2}\E[G_h^r]
&=I_h(\delta)+T_h(\delta),\\
I_h(\delta)&=h^{(r-1)K/2}\int_{B(y,\delta)}g_h(y\mid z)^r\,r_h(z\mid x)\,dz,\\
T_h(\delta)&=h^{(r-1)K/2}\int_{B(y,\delta)^c}g_h(y\mid z)^r\,r_h(z\mid x)\,dz .
\end{align*}

\emph{Tail term.} For $z\notin B(y,\delta)$ and $h\le\delta/(4\bar\mu+1)$, \ref{ass:mu} gives
\[
\lVert m_h(z)\rVert\ge\lVert y-z\rVert-h\bar\mu\ge\tfrac12\lVert y-z\rVert
\qquad\text{and}\qquad
\lVert m_h(z)\rVert\ge\tfrac12\delta .
\]
By \ref{ass:V} we have $V(z)^{-1}\succeq\overline v^{-1}I_K$ and $|V(z)|^{-r/2}\le\underline v^{-rK/2}$, so
\[
g_h(y\mid z)^r
\le(2\pi h)^{-rK/2}\underline v^{-rK/2}
\exp\!\left[-\frac{r\lVert m_h(z)\rVert^2}{2\overline vh}\right].
\]
Splitting the exponent in half and using the two lower bounds separately,
\[
g_h(y\mid z)^r
\le(2\pi h)^{-rK/2}\underline v^{-rK/2}
\exp\!\left[-\frac{r\delta^2}{16\overline vh}\right]
\exp\!\left[-\frac{r\lVert y-z\rVert^2}{16\overline vh}\right].
\]
Using \ref{ass:bound} and $\int_{\R^K}e^{-r\lVert w\rVert^2/(16\overline vh)}dw=(16\pi\overline vh/r)^{K/2}$, and collecting powers as $h^{(r-1)K/2}\cdot h^{-rK/2}\cdot h^{K/2}=1$,
\begin{equation}\label{eq:tail_bound}
T_h(\delta)\le
\bar r\,\underline v^{-rK/2}(2\pi)^{-rK/2}(16\pi\overline v/r)^{K/2}
\exp\!\left[-\frac{r\delta^2}{16\overline vh}\right]
\xrightarrow[h\downarrow0]{}0 .
\end{equation}

\emph{Local term.} Substitute $z=z_h(u)=y+\sqrt h\,u$, so that $dz=h^{K/2}du$ and the domain becomes $\lVert u\rVert<\delta/\sqrt h$. Since $h^{(r-1)K/2}(2\pi h)^{-rK/2}h^{K/2}=(2\pi)^{-rK/2}$, every power of $h$ cancels and
\[
I_h(\delta)=(2\pi)^{-rK/2}\int_{\R^K}\mathbf 1\{\lVert u\rVert<\delta/\sqrt h\}\,\Psi_h(u)\,du,
\]
\[
\Psi_h(u)=|V(z_h(u))|^{-r/2}r_h(z_h(u)\mid x)
\exp\!\left[-\frac{r}{2h}m_h(z_h(u))^{\mathsf T}V(z_h(u))^{-1}m_h(z_h(u))\right].
\]
Here $m_h(z_h(u))=-\sqrt h\,u-h\mu(z_h(u))$, so
\[
\frac1hm_h(z_h(u))^{\mathsf T}V(z_h(u))^{-1}m_h(z_h(u))
=\bigl(u+\sqrt h\,\mu(z_h(u))\bigr)^{\mathsf T}V(z_h(u))^{-1}\bigl(u+\sqrt h\,\mu(z_h(u))\bigr).
\]
Fix $u$. Then $z_h(u)\to y$, so continuity of $\mu$ and $V$ gives $V(z_h(u))\to V(y)$ and $\sqrt h\,\mu(z_h(u))\to0$, and the quadratic form converges to $u^{\mathsf T}V(y)^{-1}u$. For $h$ small enough that $\sqrt h\lVert u\rVert\le\delta_0$, \ref{ass:conv} and continuity of $p(\cdot\mid x)$ give $r_h(z_h(u)\mid x)\to p(y\mid x)$. Hence, pointwise in $u$,
\[
\Psi_h(u)\longrightarrow|V(y)|^{-r/2}p(y\mid x)\exp\!\left[-\tfrac r2u^{\mathsf T}V(y)^{-1}u\right].
\]

\emph{Domination.} Let $h\le\min\{\delta_0,(2\bar\mu)^{-2}\}$. If $\lVert u\rVert\ge2\sqrt h\,\bar\mu$ then $\lVert u+\sqrt h\,\mu(z_h(u))\rVert\ge\tfrac12\lVert u\rVert$, so by \ref{ass:V} and \ref{ass:bound}
\[
\Psi_h(u)\le\underline v^{-rK/2}\,\bar r\,\exp\!\left[-\frac{r\lVert u\rVert^2}{8\overline v}\right],
\]
while if $\lVert u\rVert<2\sqrt h\,\bar\mu\le1$ then $\Psi_h(u)\le\underline v^{-rK/2}\bar r$. The dominating function
\[
\Lambda(u)=\underline v^{-rK/2}\bar r\left\{\exp\!\left[-\frac{r\lVert u\rVert^2}{8\overline v}\right]+\mathbf 1\{\lVert u\rVert\le1\}\right\}
\]
is integrable and does not depend on $h$. Dominated convergence therefore gives
\[
\lim_{h\downarrow0}I_h(\delta)
=(2\pi)^{-rK/2}|V(y)|^{-r/2}p(y\mid x)
\int_{\R^K}\exp\!\left[-\tfrac r2u^{\mathsf T}V(y)^{-1}u\right]du .
\]
The Gaussian integral equals $(2\pi/r)^{K/2}|V(y)|^{1/2}$, so
\[
\lim_{h\downarrow0}I_h(\delta)
=(2\pi)^{-rK/2}(2\pi/r)^{K/2}p(y\mid x)|V(y)|^{-(r-1)/2}
=c_{r,K}\,p(y\mid x)|V(y)|^{-(r-1)/2}.
\]
This limit does not depend on $\delta$, and $T_h(\delta)\to0$ by \eqref{eq:tail_bound}, which proves \eqref{eq:general_r_moment}. Taking $r=2$ and using $(2\pi)^{-K/2}2^{-K/2}=(4\pi)^{-K/2}$ gives \eqref{eq:general_second}.
\end{proof}

\begin{remark}[Verification in exactly solvable cases]\label{rem:verification}
The constant can be checked against closed forms. For standard Brownian motion $V\equiv I_K$, and substituting $p(y\mid x)=(2\pi)^{-K/2}e^{-\lVert y-x\rVert^2/2}$ into \eqref{eq:general_r_moment} reproduces the exact limit of \Cref{prop:brownian_moments} term by term. For a constant-coefficient diffusion $dY=\mu\,dt+\Sigma\,dW$ the same computation with $V=\Sigma\Sigma^{\mathsf T}$ reproduces \eqref{eq:general_r_moment} including the determinant factor. For a multidimensional Ornstein--Uhlenbeck process the preterminal law is Gaussian and the second moment is again available in closed form; \Cref{fig:ou_scaling} confirms convergence to \eqref{eq:general_second} numerically. The cases $r=2$ and $r=3$ supply exactly the second and third moments used in \Cref{thm:generic_clt}.
\end{remark}

\begin{remark}[Localisation]\label{rem:localisation}
Only three features of \Cref{ass:local} are used: ellipticity and boundedness of $\mu$ near $y$, which drive the local term; the uniform bound \ref{ass:bound}, which drives domination; and enough global control to make the tail negligible. Global uniform ellipticity is therefore stronger than necessary. It can be replaced by ellipticity on $B(y,\delta_0)$ together with any condition forcing the contribution of $\{Z_h\notin B(y,\delta)\}$ to vanish after multiplication by $h^{(r-1)K/2}\sup_zg_h(y\mid z)^r$, for instance a Gaussian upper bound of Aronson type on $r_h$. We state the global version because it is directly verifiable and because the content of the theorem is the dimensional constant, not the weakest hypotheses.
\end{remark}

\begin{remark}[Time-inhomogeneous coefficients]\label{rem:time_inhomogeneous}
If $\mu$ and $V$ depend on time, the one-step kernel \eqref{eq:one_step_kernel} evaluates them at time $1-h$, and the proof goes through verbatim provided $\mu$ and $V$ are jointly continuous at $(y,1)$ and \ref{ass:mu} and \ref{ass:V} hold uniformly in $t$ on a neighbourhood of $1$. The limit constant is unchanged and $V(y)$ is read as $V(y,1)$. We suppress the time argument in the statement because carrying it obscures the cancellation of powers of $h$, which is where the dimensional constant comes from.
\end{remark}

\begin{corollary}[Generic variance scale]\label{cor:generic_variance}
Let $q_M(y\mid x)=\E[G_h]$ and suppose $q_M(y\mid x)\to p(y\mid x)$. Then
\begin{equation}\label{eq:generic_variance}
\Var(\widehat q_{M,S})
\sim
\frac{\tau^2(x,y)}{S h^{K/2}},
\qquad
\tau^2(x,y)=(4\pi)^{-K/2}p(y\mid x)|V(y)|^{-1/2}.
\end{equation}
\end{corollary}

\begin{theorem}[Generic endpoint CLT]\label{thm:generic_clt}
Let $(M_n,S_n)$ satisfy $M_n\to\infty$ and $S_n\to\infty$, write $h_n=1/M_n$, and suppose the effective endpoint size diverges:
\begin{equation}\label{eq:effective_diverge_general}
R_n=S_nh_n^{K/2}=\frac{S_n}{M_n^{K/2}}\longrightarrow\infty .
\end{equation}
Then under \Cref{ass:local},
\begin{equation}\label{eq:generic_clt}
\sqrt{R_n}
\left\{\widehat q_{M_n,S_n}(y\mid x)-q_{M_n}(y\mid x)\right\}
\Longrightarrow
\Normal\!\left(0,\tau^2(x,y)\right),
\end{equation}
where, as in \eqref{eq:generic_variance},
\begin{equation}\label{eq:tau_restated}
\tau^2(x,y)=(4\pi)^{-K/2}p(y\mid x)|V(y)|^{-1/2}.
\end{equation}
This is a statement about the simulated transition density, centred at the Euler density $q_{M_n}$ and not at the diffusion density $p$; \Cref{cor:diffusion_clt} treats the latter. It is not a statement about any parameter estimator.
\end{theorem}

\begin{proof}
Define the triangular array
\[
G_{n,s}=g_{h_n}(y\mid Z_{h_n,s}),
\qquad s=1,\ldots,S_n,
\]
where $Z_{h_n,1},\ldots,Z_{h_n,S_n}$ are the independent preterminal draws in \eqref{eq:pedersen_estimator}. Within row $n$ these are independent and identically distributed with common mean $\E G_{n,1}=q_{M_n}(y\mid x)$, and
\[
\widehat q_{M_n,S_n}-q_{M_n}=\frac1{S_n}\sum_{s=1}^{S_n}\left(G_{n,s}-\E G_{n,s}\right).
\]

\emph{Variance.} \Cref{thm:local_moments} with $r=2$ gives $h_n^{K/2}\E[G_{n,1}^2]\to\tau^2$. Since $\E G_{n,1}=q_{M_n}\to p(y\mid x)$ is bounded, $h_n^{K/2}(\E G_{n,1})^2\to0$, so
\begin{equation}\label{eq:var_asymptotic}
h_n^{K/2}\Var(G_{n,1})\longrightarrow\tau^2,
\qquad
\Var\bigl(\widehat q_{M_n,S_n}\bigr)=\frac{\Var(G_{n,1})}{S_n}\sim\frac{\tau^2}{R_n}.
\end{equation}

\emph{Third centred moment.} \Cref{thm:local_moments} with $r=3$ bounds the \emph{raw} moment, $\E[G_{n,1}^3]=O(h_n^{-K})$. The centred moment does not follow from this automatically, and is obtained as follows. The summand is a Gaussian density value, hence nonnegative, so the elementary inequality $|a-b|^3\le4(a^3+b^3)$ for $a,b\ge0$ applies and
\begin{equation}\label{eq:third_centred}
\E\lvert G_{n,1}-\E G_{n,1}\rvert^3
\le4\left\{\E[G_{n,1}^3]+(\E G_{n,1})^3\right\}
=O(h_n^{-K})+O(1)=O(h_n^{-K}),
\end{equation}
using once more that $q_{M_n}$ is bounded. Nonnegativity of the summand is what makes this step elementary; without it the centred third moment would require a separate argument.

\emph{Lyapunov condition.} Write $\sigma_n^2=\Var(G_{n,1})$, so $\sigma_n^2\sim\tau^2h_n^{-K/2}$ by \eqref{eq:var_asymptotic}. The Lyapunov ratio at order $2+\delta$ with $\delta=1$ is
\[
L_n=\frac{S_n\,\E\lvert G_{n,1}-\E G_{n,1}\rvert^3}{\bigl(S_n\sigma_n^2\bigr)^{3/2}}
=\frac{S_n\,O(h_n^{-K})}{S_n^{3/2}\,\tau^3h_n^{-3K/4}\{1+o(1)\}}
=O\!\left(S_n^{-1/2}h_n^{-K/4}\right)
=O\!\left(R_n^{-1/2}\right),
\]
because $S_n^{-1/2}h_n^{-K/4}=(S_nh_n^{K/2})^{-1/2}$. Condition \eqref{eq:effective_diverge_general} gives $L_n\to0$.

The Lyapunov central limit theorem for triangular arrays with row-wise independent summands therefore applies to $(S_n\sigma_n^2)^{-1/2}\sum_s(G_{n,s}-\E G_{n,s})$. Multiplying by $\sigma_nh_n^{K/4}\to\tau$ converts the normalisation $\sqrt{S_n}\,\sigma_n$ into $\sqrt{R_n}$ and yields \eqref{eq:generic_clt}.
\end{proof}

\begin{remark}[Why a fixed-$M$ central limit theorem cannot be reused]\label{rem:not_uniform}
For fixed $M$ the summands do not change with $S$, their variance is a constant, and the classical $\sqrt S$ central limit theorem applies. The array above is genuinely triangular: the law of $G_{n,s}$ changes with $n$ through $h_n$, and its variance diverges at rate $h_n^{-K/2}$. The correct normalisation is $\sqrt{R_n}$, which coincides with $\sqrt{S_n}$ only when $K=0$. Substituting a fixed-$M$ central limit theorem and then letting $M$ grow, as in Lemma 3 of \citet{BrandtSantaClara2002}, is exactly the step that fails; the limiting variance stated there is the variance of a summand that itself depends on $M$.
\end{remark}

\begin{corollary}[CLT centred at the diffusion density]\label{cor:diffusion_clt}
Suppose in addition that
\begin{equation}\label{eq:euler_bias_expansion}
q_M(y\mid x)-p(y\mid x)=b(x,y)h+o(h).
\end{equation}
If
\begin{equation}\label{eq:bias_negligible}
S h^{K/2}\to\infty,
\qquad
\sqrt{S h^{K/2}}\,h\to0,
\end{equation}
then
\begin{equation}\label{eq:diffusion_clt}
\sqrt{S h^{K/2}}
\left\{\widehat q_{M,S}(y\mid x)-p(y\mid x)\right\}
\Longrightarrow
\Normal(0,\tau^2(x,y)).
\end{equation}
In terms of $M$ and $S$, the second condition in \eqref{eq:bias_negligible} is
\begin{equation}\label{eq:correct_bias_rate}
\frac{\sqrt S}{M^{1+K/4}}\longrightarrow0.
\end{equation}
For $K=4$, the two conditions together define the feasible regime
\begin{equation}\label{eq:k4_feasible}
M^2\ll S\ll M^4.
\end{equation}
\end{corollary}

\begin{proof}
Decompose
\[
\sqrt{Sh^{K/2}}\left\{\widehat q_{M,S}-p\right\}
=\sqrt{Sh^{K/2}}\left\{\widehat q_{M,S}-q_M\right\}
+\sqrt{Sh^{K/2}}\left\{q_M-p\right\}.
\]
The first term converges to $\Normal(0,\tau^2)$ by \Cref{thm:generic_clt}. By \eqref{eq:euler_bias_expansion} the second term equals
\[
\sqrt{Sh^{K/2}}\left\{b(x,y)h+o(h)\right\},
\]
which vanishes precisely under the second condition in \eqref{eq:bias_negligible}. Slutsky's theorem gives \eqref{eq:diffusion_clt}.

For the restatement, $\sqrt{Sh^{K/2}}\,h=\sqrt S\,h^{K/4+1}=\sqrt S/M^{1+K/4}$, since $h=1/M$. This is \eqref{eq:correct_bias_rate}, and squaring gives the equivalent form $S/M^{2+K/2}\to0$.

For $K=4$ the first condition in \eqref{eq:bias_negligible} reads $S/M^2\to\infty$ and the second reads $S/M^{4}\to0$, which is \eqref{eq:k4_feasible}. The regime is nonempty: $S=M^3$ satisfies both.
\end{proof}

\begin{remark}[Bias negligibility is not variance consistency]\label{rem:bias_vs_variance}
The two conditions in \eqref{eq:bias_negligible} do different work and pull in opposite directions. The first is $S/M^{K/2}\to\infty$: it makes the Monte Carlo error vanish and is what \Cref{thm:generic_clt} needs. It is a \emph{lower} bound on $S$, requiring $S\gg M^{K/2}$. The second is $S/M^{2+K/2}\to0$, equivalently $\sqrt S/M^{1+K/4}\to0$: it makes the Euler bias negligible \emph{relative to the Monte Carlo standard deviation}, and because that standard deviation shrinks as $S$ grows, it is an \emph{upper} bound on $S$, requiring $S\ll M^{2+K/2}$. At $K=4$ the pair is $M^2\ll S\ll M^4$. Conflating the two is the source of the difficulty in \Cref{tab:rates}: the published condition $\sqrt S/M\to0$, that is $S\ll M^2$, has the form of the second requirement but at the wrong power, and it then contradicts the first when $K\ge4$.
\end{remark}

\begin{remark}[Dependency structure]\label{rem:dependencies}
The density-level results depend on one another as follows. \Cref{ass:local} supports \Cref{thm:local_moments}, which is used only through its cases $r=2$ and $r=3$. The case $r=2$ gives the variance scale \Cref{cor:generic_variance}; the two cases together give the Lyapunov condition in \Cref{thm:generic_clt}. \Cref{cor:diffusion_clt} adds the Euler bias expansion \eqref{eq:euler_bias_expansion}, which is an assumption imported from \citet{BallyTalay1996a,BallyTalay1996b} rather than proved here, and combines it with \Cref{thm:generic_clt}. \Cref{prop:optimal_m} uses \Cref{cor:generic_variance} and \eqref{eq:euler_bias_expansion} but no central limit theorem. In the Brownian benchmark, \Cref{prop:brownian_moments} replaces \Cref{thm:local_moments} with an exact computation and the bias expansion is unnecessary because the bias is zero, so \Cref{thm:brownian_clt} rests on no imported result at all. \Cref{thm:collapse} is independent of every central limit theorem above; it uses only the exact Gaussian tail bound.
\end{remark}

\begin{remark}[The dimensional conflict]
The condition $\sqrt S/M\to0$ suppresses an $O(M^{-1})$ Euler bias under an incorrectly fixed Monte Carlo scale. Once the variance inflation $M^{K/2}$ is included, the bias must be compared with the standard deviation $M^{K/4}/\sqrt S$, producing \eqref{eq:correct_bias_rate}. In four dimensions, the lower consistency requirement and the upper bias requirement leave the nonempty interval \eqref{eq:k4_feasible}. The published condition places $S$ below the lower boundary of that interval.
\end{remark}

\section{Bias--variance trade-off and dimensional rate}\label{sec:mse}

Suppose the Euler density admits \eqref{eq:euler_bias_expansion}. Because \eqref{eq:pedersen_estimator} is unbiased for $q_M$,
\begin{align}
\E\!\left[(\widehat q_{M,S}-p)^2\right]
&=(q_M-p)^2+\Var(\widehat q_{M,S}) \nonumber\\
&=b(x,y)^2h^2
+\frac{\tau^2(x,y)}{S h^{K/2}}
+o\!\left(h^2+\frac{1}{S h^{K/2}}\right). \label{eq:mse_expansion}
\end{align}

\begin{proposition}[Pointwise optimal discretisation]\label{prop:optimal_m}
If $b(x,y)\ne0$, the leading terms in \eqref{eq:mse_expansion} are balanced by
\begin{equation}\label{eq:optimal_h}
h\asymp S^{-2/(K+4)},
\qquad
M\asymp S^{2/(K+4)}.
\end{equation}
The resulting optimal pointwise mean squared error is
\begin{equation}\label{eq:optimal_mse}
\operatorname{MSE}_{\mathrm{opt}}
\asymp S^{-4/(K+4)}.
\end{equation}
\end{proposition}

\begin{proof}
Differentiate the leading expression $B h^2+C/(S h^{K/2})$ with respect to $h$ or equate its two terms. This gives $h^{K/2+2}\asymp S^{-1}$ and hence \eqref{eq:optimal_h}. Substitution yields \eqref{eq:optimal_mse}.
\end{proof}

The rate is the familiar structure of a kernel-density problem: reducing the bandwidth controls approximation bias but decreases the effective number of simulations. The dimension enters through the volume of the final Gaussian kernel. Some illustrative exponents are
\begin{center}
\begin{tabular}{ccc}
\toprule
Dimension $K$ & $M_{\mathrm{opt}}$ & Optimal MSE \\
\midrule
$1$ & $S^{2/5}$ & $S^{-4/5}$ \\
$2$ & $S^{1/3}$ & $S^{-2/3}$ \\
$4$ & $S^{1/4}$ & $S^{-1/2}$ \\
$8$ & $S^{1/6}$ & $S^{-1/3}$ \\
\bottomrule
\end{tabular}
\end{center}

This does not mean that $M$ should always be chosen by pointwise MSE. Likelihood estimation requires accuracy over many endpoints and parameter values, and higher-order or bridge-based methods can change both constants and rates. It does mean that increasing $M$ while holding $S$ fixed, or while letting $S$ grow too slowly relative to dimension, can make the simulation approximation worse rather than better.

\begin{table}[t]
\centering
\caption{Principal endpoint-density rate conditions.}
\label{tab:rates}
\begin{tabular}{lll}
\toprule
Property & General dimension $K$ & Four dimensions \\
\midrule
Monte Carlo mean-square consistency
& $S/M^{K/2}\to\infty$ & $S/M^2\to\infty$ \\
Density CLT centred at $q_M$
& $S/M^{K/2}\to\infty$ & $S/M^2\to\infty$ \\
Euler bias negligible in corrected CLT
& $\sqrt S/M^{1+K/4}\to0$ & $S/M^4\to0$ \\
Condition in \citet{BrandtSantaClara2002}
& $\sqrt S/M\to0$ & $S/M^2\to0$ \\
\bottomrule
\end{tabular}
\end{table}

\section{From density error to log-likelihood error}\label{sec:loglik}

For observations $Y_0,\ldots,Y_N$, the simulated log likelihood is a sum of terms
\[
\log\widehat q_{M,S}(Y_{n+1}\mid Y_n;\theta).
\]
Even when $\widehat q_{M,S}$ is unbiased for $q_M$, its logarithm is not unbiased. The effective endpoint size $R_{M,S}=S h^{K/2}$ controls this nonlinearity.

The next result is conditional. Its first two statements follow from \Cref{thm:generic_clt}, but the third holds only under an additional uniform-integrability and lower-tail condition, stated in the proposition and assumed rather than verified.

\begin{proposition}[Log-density delta expansion]\label{prop:log_delta}
Fix an endpoint with $q_M\to p>0$ and suppose the conditions of \Cref{thm:generic_clt} hold. If $R_{M,S}=S h^{K/2}\to\infty$, then
\begin{equation}\label{eq:log_clt}
\sqrt{R_{M,S}}
\left\{\log\widehat q_{M,S}-\log q_M\right\}
\Longrightarrow
\Normal\!\left(0,\frac{\tau^2}{p^2}\right),
\end{equation}
and
\begin{equation}\label{eq:log_stochastic_expansion}
\log\widehat q_{M,S}-\log q_M
=
\frac{\widehat q_{M,S}-q_M}{q_M}
-\frac{(\widehat q_{M,S}-q_M)^2}{2q_M^2}
+\op(R_{M,S}^{-1}).
\end{equation}
Whenever the second-order expansion is uniformly integrable, it follows that
\begin{equation}\label{eq:log_bias}
\E[\log\widehat q_{M,S}]-\log q_M
=-\frac{\tau^2}{2p^2R_{M,S}}+o(R_{M,S}^{-1}).
\end{equation}
\end{proposition}

\begin{proof}
The first statement is the delta method applied to \Cref{thm:generic_clt}. Since $(\widehat q-q_M)/q_M=\Op(R^{-1/2})$ and $q_M$ is bounded away from zero, Taylor expansion of $\log(1+u)$ gives \eqref{eq:log_stochastic_expansion}. Taking expectations requires uniform integrability of the remainder and lower-tail control; under that additional condition, the first-order term has zero expectation and the second-order term gives \eqref{eq:log_bias}.
\end{proof}

The leading nonlinear scale is
\begin{equation}\label{eq:log_bias_order}
R_{M,S}^{-1}=\frac{M^{K/2}}{S},
\end{equation}
not uniformly $1/S$. The statement in \citet[p.~171]{BrandtSantaClara2002} that the logarithmic transformation induces a bias of order $1/S$ is valid with $M$ fixed. It is not a joint-$M,S$ order.

Suppose simulations are independent across observations and the local constants are uniformly bounded. For the average simulated criterion
\[
\widehat Q_{N,M,S}(\theta)
=\frac1N\sum_{n=0}^{N-1}
\log\widehat q_{n,M,S}(\theta),
\]
the stochastic average of the first-order simulation error has standard deviation of order
\[
(NR_{M,S})^{-1/2},
\]
while the average nonlinear bias remains of order $R_{M,S}^{-1}$. The simulation bias does not average out across observations.

A full estimator theorem requires uniform control in $\theta$ of densities, scores, and Hessians. The following statement is an abstract perturbation result: it is a deterministic implication about any pair of maximisers satisfying the listed conditions, and it is \emph{not} a theorem about the simulated-likelihood estimator, because we do not verify conditions (ii) and (iii) for that estimator. It is included to separate what the endpoint calculations deliver from what an estimator theorem would additionally require.

\begin{proposition}[Abstract equivalence to exact MLE]\label{prop:m_estimator}
Let $\widehat\theta_N$ maximise an exact average log likelihood $Q_N$, and let $\widehat\theta_{N,M,S}$ maximise a simulated criterion $\widehat Q_{N,M,S}$. Assume:
\begin{enumerate}[label=(\roman*)]
\item $\widehat\theta_N\to\thetao$ and $\sqrt N(\widehat\theta_N-\thetao)$ has the usual exact-MLE limit;
\item on a neighbourhood of $\thetao$,
\[
\sup_\theta
\|\nabla\widehat Q_{N,M,S}(\theta)-\nabla Q_N(\theta)\|
=o_{\Pp}(N^{-1/2});
\]
\item on the same neighbourhood,
\[
\sup_\theta
\|\nabla^2\widehat Q_{N,M,S}(\theta)-\nabla^2Q_N(\theta)\|
=o_{\Pp}(1);
\]
\item the exact criterion has a nonsingular limiting Hessian.
\end{enumerate}
Then
\[
\sqrt N(\widehat\theta_{N,M,S}-\widehat\theta_N)\to0
\quad\text{in probability}.
\]
\end{proposition}

\begin{proof}
Expand the simulated first-order condition around $\widehat\theta_N$ and use the uniform score and Hessian bounds together with nonsingularity of the limiting exact Hessian.
\end{proof}

Under uniform analogues of the pointwise density expansion, plausible minimum requirements for first-order equivalence include
\begin{equation}\label{eq:parameter_rate_heuristic}
S h^{K/2}\to\infty,
\qquad
\frac{\sqrt N}{S h^{K/2}}\to0,
\qquad
\sqrt N\,h\to0,
\end{equation}
corresponding respectively to vanishing score simulation noise, vanishing nonlinear simulation bias at the $N^{-1/2}$ scale, and negligible Euler bias. Derivatives of the simulated paths and Gaussian kernels can impose stronger conditions. Therefore \eqref{eq:parameter_rate_heuristic} is a design implication, not a claimed complete theorem for arbitrary diffusions.

The broader Monte Carlo likelihood literature treats joint data and simulation limits through empirical-process, importance-sampling, or exact-simulation arguments rather than by combining two pointwise limits; see, among others, \citet{Lee1995,SungGeyer2007,BeskosEtAl2009,MiasojedowEtAl2016}.

\section{What fails in the published asymptotic argument}\label{sec:proof_review}

This section distinguishes statements refuted by the Brownian example from steps that are merely unproved under the stated assumptions. Two published statements are directly refuted: the joint transition-density consistency of Lemma 2 and the joint central-limit statement of Lemma 3, by \Cref{thm:collapse} and \Cref{cor:clt_failure} respectively. Everything else below is a gap in a proof, not a demonstrated falsehood. \Cref{rem:hierarchy} states the hierarchy compactly.

\subsection{Lemma 2: sequential convergence is used as joint convergence}

The proof of Lemma 2 in \citet{BrandtSantaClara2002} applies the strong law to the Monte Carlo average and then invokes convergence of the Euler density. For every fixed $M$,
\[
\widehat q_{M,S}\to q_M\quad\text{as }S\to\infty,
\]
and separately
\[
q_M\to p\quad\text{as }M\to\infty.
\]
This gives the iterated limit
\[
\lim_{M\to\infty}\lim_{S\to\infty}\widehat q_{M,S}=p.
\]
It does not give convergence along every sequence $(M_n,S_n)\to(\infty,\infty)$. The summand changes with $M$ and its variance diverges as $M^{K/2}$. \Cref{thm:collapse} directly refutes the stated joint conclusion.

Two features of Lemma 2 are worth isolating. It is stated as a joint limit in $M$ and $S$ with no rate restriction whatsoever, so the refutation does not depend on the separate discussion of $\sqrt S/M\to0$ below; and it is stated under Assumptions 1 and 2 alone, which standard Brownian motion satisfies. The counterexample therefore meets the lemma on its own terms.

Because Lemmas 3 to 6 take Lemma 2 as an input, and Theorem 1 is stated as the summary of Lemmas 1 to 6, the published proofs of Theorems 1 and 2 fail. That is a statement about the proofs. It does not establish that the conclusions of those theorems are false, and the subsections that follow identify further steps that are independently unproved.

\subsection{Lemma 3: the joint central-limit statement is directly false}

The paper uses a $\sqrt S$ central limit theorem and requires $\sqrt S/M\to0$ to suppress the Euler bias. Two objections apply, and they are of different strengths.

The weaker objection is that the statement is ill-posed. The variance of the triangular-array summand is not asymptotically constant; it is of order $M^{K/2}$, so the limiting variance written in Lemma 3, namely the variance of $\phi(Y_{t_{n+1}};z_s+\mu h,Vh)$, is the variance of a quantity that itself depends on $M$ and diverges. The correct standardisation is $\sqrt{S/M^{K/2}}$, and the Euler bias must then be compared with that scale, producing \eqref{eq:correct_bias_rate}.

The stronger objection is that the statement is false, not merely ill-posed. \Cref{cor:clt_failure} shows that along $M_n=S_n=n$, a sequence satisfying the stated rate condition, $\sqrt{S_n}(\widehat q_{n,n}-p)\to-\infty$ in probability. The sequence is not tight, so it converges in distribution to no limit at all, and no choice of limiting variance rescues the statement. This refutation is independent of the objection about ill-posedness and does not rely on identifying what the intended limiting variance was.

\subsection{Lemma 4: an \texorpdfstring{$O_{\Pp}$}{O\_P} term is treated as \texorpdfstring{$o_{\Pp}$}{o\_P}}

Let $x_n=\widehat q_n-p_n$. A nondegenerate central limit theorem implies $x_n=O_{\Pp}(S^{-1/2})$ for fixed $M$, not $o_{\Pp}(S^{-1/2})$. The Taylor expansion
\[
\log\left(1+\frac{x_n}{p_n}\right)
=\frac{x_n}{p_n}+O_{\Pp}(x_n^2)
\]
retains a first-order stochastic term. In the joint limit, the correct order is $O_{\Pp}\{(S h^{K/2})^{-1/2}\}$. Moreover, summing pointwise errors over $N$ observations requires lower bounds or tail control for $p_n$, dependence control, and a triangular-array argument. These are not supplied.

\subsection{Lemma 5: convergence of objective values does not control maximisers}

An argmax result requires uniform convergence of the criterion and separation or local curvature around the maximiser. A bounded Hessian gives an upper quadratic bound. To infer that two maximisers are close from a small objective difference, one needs a lower quadratic bound or a negative Hessian whose smallest eigenvalue is bounded away from zero after normalisation. Pointwise convergence and identification alone are insufficient.

The structure of the argument makes this explicit. Summing the two second-order expansions leaves a difference of objective values on the left and a term proportional to $(\widehat\theta_N-\widehat\theta_{N,M,S})^2$ on the right, whose coefficient is an average of two Hessians. Assumption 4 of \citet{BrandtSantaClara2002} bounds that coefficient above. Dividing a small left-hand side by a coefficient that is only bounded above yields no bound at all on the squared distance; what is required is a bound below, that is, a uniformly negative definite Hessian in a neighbourhood. The conclusion drawn, $\widehat\theta_{N,M,S}-\widehat\theta_N=o(N^{1/2}/S^{1/4})$, also inherits whatever is wrong with the $o(N/S^{1/2})$ input from Lemma 4.

\subsection{Lemmas 6 and 7: exact-MLE asymptotics require additional structure}

Consistency of the exact diffusion MLE ordinarily requires a uniform law of large numbers, global or local identification, and stability assumptions such as stationarity and ergodicity when $N\to\infty$. A score expansion around $\thetao$ is not by itself a proof of consistency because the intermediate Hessian is evaluated between the estimator and $\thetao$.

Similarly, the martingale central limit theorem invoked from \citet{HallHeyde1980} requires convergence of the realised conditional quadratic variation and a conditional Lindeberg condition. Defining Fisher information as an expectation does not make the realised conditional variance condition automatic. Boundedness of $\partial p/\partial\theta$ is also not a bound on the score $(\partial p/\partial\theta)/p$ when $p$ can be small; strict positivity of $p$ at each point is weaker than the uniform lower bound the argument needs.

These points do not imply that no correct exact-MLE theorem can be written. They imply that the short proof in Appendix A does not establish the stated result, even apart from the false transition-density joint limit.

\subsection{What a counterexample to the argmax estimator would require}\label{sec:argmax_gap}

It is worth being explicit about the limit of the negative results above. \Cref{thm:collapse} refutes a statement about the simulated transition density at one pair of points. It does not refute the conclusion of Theorem 1, which concerns the maximiser of the simulated log likelihood. Failure of a criterion pointwise does not imply failure of its argmax, and three obstacles separate the two.

First, the collapse is driven by a rare-event geometry that is not uniform over the state space or over the parameter. \Cref{thm:collapse} fixes $x=y=0$. An argmax argument needs control of $\widehat q_{M,S}(Y_{n+1}\mid Y_n;\theta)$ simultaneously over $N$ observed pairs and over a neighbourhood of $\thetao$, and the sets on which individual densities collapse need not align.

Second, the simulated criterion is evaluated with common random numbers across $\theta$, as the published implementation requires for a smooth objective. The summands are therefore strongly dependent across parameter values. A criterion that is badly behaved at each fixed $\theta$ may still have a well-behaved maximiser if the errors move together; conversely, common random numbers can create a random limit criterion whose argmax is random. Which happens is a property of the coupling, not of the pointwise limit.

Third, the logarithm is unbounded below. Collapse towards zero at a single observation drives one log-density term to $-\infty$, which may push the maximiser away from the region where collapse occurs rather than towards a wrong value. The direction of that effect depends on how the collapse region varies with $\theta$.

A complete argmax counterexample would need to exhibit a parametric family, an admissible sequence $(M_n,S_n)$, and either a limit criterion whose maximiser differs from $\thetao$, or a demonstration that the simulated criterion fails to converge uniformly on a neighbourhood in a way that the maximiser inherits. We record in \Cref{app:argmax} the approaches we attempted and where each fails. Until such a construction exists, the correct statement is that the published proofs of Theorems 1 and 2 do not establish their conclusions, not that those conclusions are false.

\section{Mathematical status of the exchange-rate application}\label{sec:application}

The empirical application estimates a four-dimensional system consisting of two interest rates, the log exchange rate, and exchange-rate volatility. It also reconstructs a market-incompleteness state. This section states the findings that bear on one question: whether the generic theorem of \Cref{sec:general} covers the model that was actually estimated. The supporting analysis is in \Cref{app:sqrt} to \Cref{app:constrained}, and every claim below is proved there.

None of what follows is an asymptotic theorem. These are calculations about one published model, reported as diagnostics.

\subsection{Summary of findings}\label{sec:application_summary}

\begin{enumerate}[label=(\roman*)]
\item \emph{The coefficients are outside the assumptions, but not equally so.} The interest rates follow square-root processes, which violate \ref{ass:V} at the origin. At the reported estimates this is a gap in proof coverage rather than a numerical hazard: the Feller ratios range from $6.4$ to $48.5$, so the origin is unattainable, and although the one-step Euler negativity probabilities are strictly positive they are of order $10^{-2667}$ or smaller at the long-run means. The incompleteness state is different. Its Feller ratio is exactly $1/2$ for every admissible parameter value, so no estimate can move it away from its boundary, and the worst-case one-step Euler negativity probability equals $\Phi(-\sqrt{1-2\alpha h})$ exactly, which tends to $\Phi(-1)\approx0.16$ as the step is refined rather than falling. See \Cref{app:sqrt}.
\item \emph{The estimator does not simulate the state whose boundary is at issue.} The implementation eliminates the incompleteness state and simulates $(r,r^*,e,v)$, recovering $\epsilon_t$ from a square root of the variance identity. Every Euler substep therefore requires that radicand to stay nonnegative at a \emph{simulated} state, not merely at the observed data, and a violation makes the next substep's coefficients complex rather than merely inaccurate. See \Cref{app:reality}.
\item \emph{The two state representations are not interchangeable.} Taking volatility as primitive, which is what the estimator does, the variance identity reconstructs incompleteness uniquely. Taking incompleteness as primitive, recovering positive volatility requires solving a quadratic, and at the reported U.S.--U.K. loadings the inverse map is not globally single-valued: the quadratic has either no positive root or two. A branch-selection condition would be needed to read the specification that way, and none is supplied. See \Cref{app:invert}.
\item \emph{Correlation validity holds wherever we can check it; only the global guarantee is missing.} The reported matrices are positive definite at every model-feasible reference state examined, with minimum eigenvalues from $0.68$ down to $0.013$ at rates a hundred times their long-run means, and no violation arose in $5{,}000$ perturbation draws per system. We find no model-feasible state at which the matrix is indefinite, in any of the three specifications, once the time-varying cases are evaluated by solving the variance identity for volatility rather than reading it from data. What remains is that the variance identity already forces every individual correlation into $[-1,1]$, so the entrywise constraint the paper imposes adds nothing, while the joint condition is nowhere enforced. A separate finding emerges from the same sweep: for the U.S.--U.K. model-C loadings, two thirds of the swept states admit no positive volatility at all, and none admits exactly one. See \Cref{app:psd}.
\item \emph{Identification is by a corner rule, not by the likelihood.} The minimum-incompleteness criterion is affine in the selected parameter, so its optimum is a vertex of the feasible set. Which constraint is active depends on the full feasible set, including range and semidefiniteness restrictions. See \Cref{app:mincomp}.
\item \emph{The final estimator is not the one the theorem describes.} It combines simulated likelihood, auxiliary least squares, a corner selection rule, and inequality constraints that the paper reports as binding. See \Cref{app:constrained}.
\end{enumerate}

Taken together these say that the estimator analysed in \Cref{sec:general} and the estimator implemented in the application are different objects, quite apart from the rate conditions. The remainder of this section reports the one diagnostic that concerns the simulation design itself.

\subsection{Finite-sample scaling diagnostics for the implementation}\label{sec:design}

Before turning to the structure of the model it is worth recording how the implemented simulation design sits relative to the asymptotic ordering that the published theorems describe. Theorem 2 of \citet{BrandtSantaClara2002} imposes two conditions jointly,
\begin{equation}\label{eq:published_joint}
\frac{\sqrt S}{M}\longrightarrow0,
\qquad
\frac{N}{S^{1/4}}\longrightarrow0,
\end{equation}
the second entering through the $o(N^{1/2}/S^{1/4})$ bound of their Lemma 5. A sequence satisfying both must eventually have $N^4\ll S\ll M^2$.

What follows is a set of diagnostics, not a test. A single finite triple $(N,M,S)$ cannot satisfy or violate a limit; only a sequence can. The question a diagnostic can answer is narrower and still useful: are the finite ratios appearing in \eqref{eq:published_joint} small, so that the implementation resembles a point along such a sequence, and does the theorem supply any error bound at the implemented sizes? The answer to the second question is no in any case, because \eqref{eq:published_joint} is a statement about limits and carries no quantitative finite-sample content.

The estimation uses $N=544$ weekly observations from January 1990 to May 2000, $M=10$ Euler substeps, and $S=5{,}000$ simulations together with $5{,}000$ antithetic variates. \Cref{tab:design} reports the ratios. At the implemented values $\sqrt S/M=7.07$ and $N/S^{1/4}=64.7$. Neither is small, so the implementation does not resemble a design drawn from a sequence along which \eqref{eq:published_joint} is operating. That is the whole of the claim; nothing here shows that the estimates are inaccurate.

The antithetic variates deserve a separate word, and a careful one. The paired innovation vectors $\varepsilon$ and $-\varepsilon$ are deterministic opposites, but the endpoint density contributions are nonlinear functions of them, so their correlation is not $-1$ and need not even be negative. It depends on the model, the state, the parameter and the endpoint. For the Brownian benchmark it can be computed exactly, and the answer is not the favourable one.

\begin{proposition}[Antithetic pairing at a coincident endpoint]\label{prop:antithetic}
For $K$-dimensional Brownian motion started at $x$, let $G$ denote the endpoint summand \eqref{eq:brownian_summand} generated from a preterminal draw $Z$ and let $G^-$ denote its antithetic partner, generated from the reflected innovations. If $y=x$ then
\begin{equation}\label{eq:antithetic_identical}
G^-=G\qquad\text{almost surely},
\end{equation}
so the pair is perfectly \emph{positively} correlated and provides no variance reduction whatsoever. An antithetic design with $S$ density evaluations then has exactly the variance of $S/2$ independent draws.
\end{proposition}

\begin{proof}
Reflecting every innovation maps the preterminal state $Z=x+\sqrt{1-h}\,\xi$ to $Z^-=x-\sqrt{1-h}\,\xi$, so $Z^--x=-(Z-x)$. The summand depends on $Z$ only through $\lVert y-Z\rVert^2$, and at $y=x$ this is $\lVert x-Z\rVert^2=\lVert Z-x\rVert^2$, which is invariant under $Z-x\mapsto-(Z-x)$. Hence $G^-=G$ pointwise. For the variance statement, averaging $S/2$ pairs each of which contributes $\tfrac12(G+G^-)=G$ gives the average of $S/2$ independent copies of $G$.
\end{proof}

More generally, writing $\rho_A$ for the correlation between a contribution and its antithetic partner, the antithetic average over $S$ evaluations has variance $(1+\rho_A)\sigma^2/S$, against $\sigma^2/S$ for $S$ independent draws. \Cref{prop:antithetic} is the case $\rho_A=+1$. Numerically, at $h=1/10$ and $K=4$, $\rho_A$ is $+1$ at $\lVert y-x\rVert=0$, $+0.063$ at $\lVert y-x\rVert=0.5$ and $-0.012$ at $\lVert y-x\rVert=1.5$; at $K=1$ and $\lVert y-x\rVert=0.5$ it is $-0.479$. Antithetic sampling therefore helps for some endpoints and hurts for others, and the effect is on the variance \emph{constant}, never on the dimensional exponent: the rate remains governed by $S/M^{K/2}$ however the pairs are correlated.

We therefore fix one convention for the whole paper. \textbf{$S$ denotes the number of independent innovation draws, equivalently the number of antithetic pairs}, and each pair yields two density evaluations. On this convention the reported design has $S=5{,}000$ independent units and $10{,}000$ density evaluations, and \Cref{tab:design} counts $S=5{,}000$. Counting the evaluations instead would overstate the number of independent units, and at coincident endpoints would overstate it by exactly a factor of two.

\begin{table}[t]
\centering
\caption{Finite-sample scaling ratios for the implemented design and three variations, at $N=544$ and $K=4$. The effective endpoint size is $R_{M,S}=S/M^{2}$ from \Cref{def:effective_size}. These are diagnostics: the columns record how the implementation is placed, not whether any limit holds.}
\label{tab:design}
\begin{tabular}{lrrrrr}
\toprule
Scenario & $M$ & $S$ & $R_{M,S}=S/M^2$ & $\sqrt S/M$ & $N/S^{1/4}$ \\
\midrule
Implemented & 10 & 5{,}000 & 50.0 & 7.07 & 64.7 \\
Double $M$ and $S$ & 20 & 10{,}000 & 25.0 & 5.00 & 54.4 \\
Double $M$, $S$ fixed & 20 & 5{,}000 & 12.5 & 3.54 & 64.7 \\
Double $M$, quadruple $S$ & 20 & 20{,}000 & 50.0 & 7.07 & 45.7 \\
\bottomrule
\end{tabular}
\end{table}

The dimensional diagnostic of \Cref{def:effective_size} gives the complementary reading. With $K=4$, $M=10$ and $S=5{,}000$,
\begin{equation}\label{eq:implemented_effective_size}
R_{M,S}=\frac{S}{M^{K/2}}=\frac{5{,}000}{100}=50 .
\end{equation}
Each transition density in a $544$-term log likelihood is therefore resolved by an effective endpoint sample of order $10^2$, not $10^4$. Again this is a scaling statement and not a verdict: an effective size of $50$ may well be adequate at these parameter values, and the exact Brownian formulas of \Cref{sec:brownian} would be the way to check that for a given target density.

The last three rows of \Cref{tab:design} record how $R_{M,S}$ responds to refinement, because the arithmetic is easy to get wrong. In four dimensions $R_{M,S}=S/M^2$, so
\begin{equation}\label{eq:doubling_arithmetic}
R_{2M,2S}=\frac{2S}{4M^2}=\frac{R_{M,S}}{2},
\qquad
R_{2M,S}=\frac{R_{M,S}}{4},
\qquad
R_{2M,4S}=R_{M,S}.
\end{equation}
The paper reports a final round of optimisation in which both $S$ and $M$ are doubled. That halves the effective endpoint size rather than increasing it, and the reported stability of the estimates under this refinement is therefore not evidence that the simulation is converging in the endpoint dimension; it is a comparison between two designs, the second of which has half the effective size of the first. Holding $R_{M,S}$ fixed while doubling $M$ requires quadrupling $S$, which is the last row of the table.

\section{Numerical evidence}\label{sec:numerics}

All numerical results in this section are generated by the accompanying Python script. They use exact Gaussian benchmarks where possible and are intended to illustrate the theorems rather than replace them.

\subsection{Second-moment scaling}

\Cref{fig:brownian_scaling} plots
\[
h^{K/2}\E[G_h^2]
\]
divided by its limiting value for $K\in\{1,2,4,8\}$. The exact formula \eqref{eq:brownian_second} converges to the constant predicted by \Cref{thm:local_moments}. The unscaled second moment grows as $M^{K/2}$.

\begin{figure}[t]
\centering
\includegraphics[width=0.86\textwidth]{figures/brownian_second_moment_scaling.pdf}
\caption{Exact local second-moment scaling for standard Brownian motion. The plotted quantity converges to one after division by its theoretical limit.}
\label{fig:brownian_scaling}
\end{figure}

\subsection{Three joint asymptotic paths}

For $K=4$, \Cref{fig:path_rmse} compares the exact relative root mean squared error under $S=M$, $S=M^2$, and $S=M^3$. The first path satisfies the published condition but has increasing RMSE. The second keeps the effective endpoint size constant and does not become consistent. The third satisfies $S/M^2\to\infty$ and has decreasing RMSE.

\begin{figure}[t]
\centering
\includegraphics[width=0.86\textwidth]{figures/brownian_joint_path_rmse.pdf}
\caption{Exact relative RMSE at the origin for four-dimensional Brownian motion.}
\label{fig:path_rmse}
\end{figure}

\subsection{Typical estimates collapse while the mean remains correct}

\Cref{fig:collapse_simulation} simulates the counterexample path $M=S$. The sample median falls rapidly towards zero, and by $M=512$ more than $95\%$ of estimates are below half the true density. The sample mean remains near one after normalisation because a small number of very large realisations carry the expectation.

\begin{figure}[t]
\centering
\includegraphics[width=0.86\textwidth]{figures/brownian_subcritical_distribution.pdf}
\caption{Monte Carlo distribution along $M=S$ for four-dimensional Brownian motion. Values are divided by the true density. The dashed line is one.}
\label{fig:collapse_simulation}
\end{figure}

\begin{table}[t]
\centering
\caption{Simulation of the four-dimensional Brownian counterexample path $M=S$. The estimator is divided by the true density.}
\label{tab:collapse_stats}
\small
\begin{tabular}{rrrrrrr}
\toprule
$M$ & Replications & Mean & Median & 10th pct. & 90th pct. & $\Pp(\widehat q<0.5p)$ \\
\midrule
8   & 20000 & 1.001 & 0.3674 & $1.427\times10^{-2}$ & 2.940 & 0.560 \\
16  & 20000 & 1.009 & 0.1582 & $1.520\times10^{-3}$ & 3.161 & 0.665 \\
32  & 20000 & 1.027 & 0.03628 & $5.925\times10^{-5}$ & 2.746 & 0.754 \\
64  & 20000 & 1.002 & 0.004223 & $8.507\times10^{-7}$ & 1.777 & 0.827 \\
128 & 20000 & 1.042 & $1.537\times10^{-4}$ & $1.487\times10^{-9}$ & 0.7409 & 0.886 \\
256 & 19531 & 0.941 & $1.118\times10^{-6}$ & $7.168\times10^{-14}$ & 0.1579 & 0.929 \\
512 & 9765  & 0.998 & $8.604\times10^{-10}$ & $7.076\times10^{-20}$ & 0.01738 & 0.955 \\
\bottomrule
\end{tabular}
\end{table}

\subsection{A nonzero drift does not remove the dimensional law}

Consider the diagonal Ornstein--Uhlenbeck process
\[
dY_t=-Y_tdt+dW_t
\]
in four dimensions. Its exact transition density is known, and the preterminal Euler state is Gaussian, so the second moment of the endpoint estimator can again be computed exactly. \Cref{fig:ou_scaling} shows convergence to the same local constant in \Cref{thm:local_moments}. The phenomenon is therefore not specific to driftless Brownian motion.

\begin{figure}[t]
\centering
\includegraphics[width=0.86\textwidth]{figures/ou_second_moment_scaling.pdf}
\caption{Scaled second moment for a four-dimensional Ornstein--Uhlenbeck diffusion.}
\label{fig:ou_scaling}
\end{figure}

\subsection{Euler boundary violations in the incompleteness process}

\Cref{fig:negative_euler} evaluates \eqref{eq:negative_prob} at the reported estimates $(\alpha,\beta)=(0.320,0.088)$ and $(0.338,0.101)$, using the step sizes $h=1/520$ and $h=1/1040$ that correspond to weekly data with $M=10$ and $M=20$, with a coarser value shown only as a labelled illustration. The horizontal scale is logarithmic so that the region near the boundary is visible; on a linear scale the entire phenomenon is compressed into a sliver next to the origin.

The figure makes \Cref{prop:worst_case_boundary} visible: refining the step translates the curve to the left without lowering it. As explained in \Cref{app:reality}, this calculation applies to the incompleteness specification \eqref{eq:epsilon_cir} treated as a state in its own right. The implemented chain simulates $(r,r^*,e,v)$ and recovers $\epsilon_t$ from \eqref{eq:epsilon_recovered}, so the boundary that must not be crossed is the complete-markets surface rather than $\{H=0\}$; the figure is an exact illustration of the mechanism rather than of the implemented chain. The contrast with the interest rates in \Cref{tab:feller} is the substantive content: for $r$ and $r^*$ the Feller ratios exceed one by an order of magnitude and the boundary is never approached, whereas for the incompleteness state the ratio is exactly $1/2$ for all parameter values, and under the conditions of \Cref{prop:min_incomplete} the boundary is reached at one sample date.

\begin{figure}[t]
\centering
\includegraphics[width=0.86\textwidth]{figures/epsilon_squared_euler_negative_probability.pdf}
\caption{One-step probability that ordinary Euler discretisation of the incompleteness specification $H_t=\epsilon_t^2$ in \eqref{eq:epsilon_cir} produces a negative state, at the estimates in Table 3 of \citet{BrandtSantaClara2002}. Solid lines are the U.S.--U.K. estimates and dashed lines the U.S.--Germany estimates. The application step sizes $h=1/520$ and $h=1/1040$ are shown together with a coarser illustrative value. Halving $h$ shifts each curve left without lowering its peak, which stays near $\Phi(-1)$; see \Cref{prop:worst_case_boundary}. The implemented estimator substitutes $\epsilon_t$ out and simulates $(r,r^*,e,v)$ instead, so this illustrates the boundary mechanism rather than the implemented chain; see \Cref{app:reality}.}
\label{fig:negative_euler}
\end{figure}

\section{Remedies and alternative constructions}\label{sec:remedies}

The endpoint variance problem is not an unavoidable property of likelihood inference for diffusions. It is a consequence of drawing preterminal states from the unconditioned forward Euler law and asking a shrinking final Gaussian kernel to connect them to a fixed observed endpoint.

\subsection{Bridge and importance proposals}

A bridge proposal guides simulated paths towards the observed endpoint. In importance-sampling language, it increases the probability of drawing paths in the region that contributes to the transition-density integral and corrects the proposal through a likelihood ratio. This is the direction taken by the acceleration methods studied by \citet{DurhamGallant2002}, the path-augmentation schemes of \citet{ElerianChibShephard2001}, conditioned-diffusion methods such as \citet{DelyonHu2006}, the regularised proposal of \citet{Lindstrom2012}, and later guided proposals for multidimensional bridges \citep{SchauerEtAl2017,VanDerMeulenSchauer2017}. The nonparametric simulation approach of \citet{Nicolau2002} replaces the endpoint Gaussian kernel by a density estimate, which changes the bandwidth problem rather than removing it.

A useful proposal should be assessed through the moments of its importance weights, and this is the respect in which bridge proposals differ in kind rather than in degree. \citet{StramerYan2007a} show that the modified Brownian bridge has importance-weight variance bounded uniformly in $M$, which is exactly what fails here: \Cref{cor:generic_variance} gives $\Var(G_h)\asymp h^{-K/2}$ for the endpoint summand. Their numerical comparison makes the contrast visible, the endpoint estimator's error increasing with $M$ at fixed $N$ while the bridge's does not, and their accuracy for the bridge is dimension-free where the endpoint rate is not. Conditioning on the endpoint therefore addresses the rare-event geometry identified here at its source, rather than merely reducing a constant.

\subsection{Exact or discretisation-free likelihood estimators}

For restricted classes of diffusions, exact simulation and unbiased transition-density estimators remove Euler bias entirely. \citet{BeskosEtAl2006,BeskosEtAl2009} develop likelihood-based methods of this kind. Their applicability is model dependent, but they illustrate the correct theoretical structure: unbiasedness, continuity in parameters, finite-moment conditions, and joint control of data and Monte Carlo sample sizes are stated explicitly.

\subsection{Analytical density expansions}

Closed-form expansions, such as those of \citet{AitSahalia2002}, avoid endpoint Monte Carlo variance but introduce an approximation-order problem of their own. They can be attractive for models that can be transformed to satisfy the required regularity conditions.

\subsection{Positivity-preserving discretisation}

For square-root components, ordinary Euler should be replaced by an exact transition where available or by a specified positivity-preserving approximation. The choice is part of the estimator, not an implementation detail, because different truncation and reflection rules define different transition densities.

\subsection{Practical design rule}

If the unconditioned endpoint estimator is nevertheless used, $M$ and $S$ should be calibrated jointly. A minimum diagnostic is the effective endpoint size
\[
R_{M,S}=S/M^{K/2}.
\]
Increasing $M$ without increasing $R_{M,S}$ is not an accuracy improvement. For density-level asymptotics, $R_{M,S}$ must diverge. For likelihood-level work, one should also monitor the variance and lower tail of every estimated transition density, use log-sum-exp evaluation, and verify stability after increasing both $M$ and $S$ along a theoretically admissible path.

\section{Conclusion}

The endpoint simulated-likelihood estimator has a simple but consequential dimensional structure. Its final Gaussian transition is a shrinking kernel. In $K$ dimensions, the second moment of one simulated contribution grows as $M^{K/2}$, so the Monte Carlo variance is of order $M^{K/2}/S$. The relevant simulation size is $S/M^{K/2}$.

This leads to four conclusions about \citet{BrandtSantaClara2002}, stated in decreasing order of strength. First, two published statements are directly false: the joint transition-density consistency of Lemma 2 and the joint central-limit statement of Lemma 3. An exact Brownian sequence satisfying the paper's own Assumptions 1 and 2, and satisfying its rate condition, makes the simulated density converge in probability to zero rather than to the positive true density, and makes $\sqrt S$ times the centred density diverge to $-\infty$. Second, because Lemmas 4 to 6 take those results as inputs, the published proofs of Theorems 1 and 2 fail. That is a statement about the proofs; we do not show that their conclusions are false, and \Cref{app:argmax} explains why the density counterexample does not transfer to the maximiser. Third, the rate condition used in those theorems is incompatible with pointwise mean-square consistency of the simulated transition density in the four-dimensional application. Fourth, and independently of all of the above, the ratios corresponding to the paper's own conditions $\sqrt S/M\to0$ and $N/S^{1/4}\to0$ are not small at the implemented design $N=544$, $M=10$, $S=5{,}000$, so that design is not in a regime resembling the asymptotic ordering the theorems describe.

The corrected transition-density analysis also supplies constructive results. The density estimator has a triangular-array CLT under $S/M^{K/2}\to\infty$. If the Euler density bias is first order, the MSE-optimal discretisation is $M\asymp S^{2/(K+4)}$. The nonlinear log-likelihood scale is $M^{K/2}/S$, subject to the usual lower-tail conditions required for expectations of logarithms. These results explain why bridge proposals and other importance-sampling methods are not merely computational refinements: they address the rare-event geometry of the estimator.

The exchange-rate application raises separate questions. Its square-root states fall outside the smooth uniformly elliptic theory, although for the interest rates the Feller ratios are large enough that this is a gap in proof coverage rather than a numerical hazard. The incompleteness state is different: its Feller ratio is exactly one half for every admissible parameter value, and under the conditions of \Cref{prop:min_incomplete} the identification rule places it on its degenerate boundary at one sample date. Because the estimator substitutes that state out and simulates volatility instead, the requirement becomes that a recovered square root remain real along every simulated Euler path, which ordinary Euler discretisation does not guarantee. The volatility equation is implicit when risk prices depend on volatility, and at the reported U.S.--U.K. estimates the inverse mapping from incompleteness to positive volatility is not globally single-valued without a branch-selection condition. The Brownian correlation matrix is positive definite everywhere we could check, including at the feasible minimum volatility; the residual point is that a joint positive-semidefinite restriction is required and entrywise bounds do not deliver it. The final empirical estimator also contains a minimum-incompleteness corner rule and binding constraints that are not covered by ordinary interior SML asymptotics.

None of these points, individually or collectively, proves that the finite-sample economic conclusions are false. They show that those conclusions cannot inherit the claimed maximum-likelihood asymptotics from the mathematical argument given. The appropriate post-pub\-li\-ca\-tion response is a corrected theorem, a clear statement of the estimator actually implemented, and a re-examination of numerical stability under dimensionally valid simulation rates or bridge-based alternatives.

\appendix

\section{Additional proof details}\label{app:proofs}

\subsection{The Brownian collapse for more general subcritical paths}

The proof of \Cref{thm:collapse} extends beyond $M=S$. Let $h_n\downarrow0$ and choose a radius
\[
r_n=\sqrt{c h_n\log(1/h_n)}.
\]
If
\[
S_n h_n^{K/2}\{\log(1/h_n)\}^{K/2}\to0
\]
and $S_n$ grows at most polynomially in $1/h_n$, then with high probability no preterminal draw lies within $r_n$ of the endpoint, while the maximum kernel outside that ball tends to zero for sufficiently large $c$. Thus the estimator again collapses to zero. The theorem uses $M=S$ because it is transparent and directly satisfies the published rate in four dimensions.

\section{Square-root coefficients violate the generic assumptions}\label{app:sqrt}

The interest rates are modelled as correlated Cox--Ingersoll--Ross processes \citep{CoxIngersollRoss1985},
\begin{equation}\label{eq:cir_rates}
dr_t=\lambda(\vartheta-r_t)dt+\sigma\sqrt{r_t}\,dW_t,
\qquad
d r_t^*=\lambda^*(\vartheta^*-r_t^*)dt+\sigma^*\sqrt{r_t^*}\,dW_t^*.
\end{equation}
The generic theory assumes infinitely differentiable coefficients with bounded derivatives and a positive-definite covariance matrix. The function $r\mapsto\sqrt r$ is not differentiable at zero, has unbounded derivatives near zero, is not defined on the negative half-line as a real coefficient, and makes the covariance degenerate at the boundary. The same singularity appears explicitly in the transformed system of Appendix B of \citet{BrandtSantaClara2002}, whose drift for $\sqrt{r_t}$ contains the term $-\tfrac18\sigma^2r_t^{-1/2}$.

It is important to separate this structural violation from its practical consequences, and the two point in opposite directions for the interest rates. \Cref{tab:feller} reports the Feller ratio $2\lambda\vartheta/\sigma^2$ at the estimates of Table 3 of \citet{BrandtSantaClara2002}. It ranges from $6.4$ to $48.5$, so in every case the ratio comfortably exceeds one, the origin is unattainable for the exact processes, and the interest-rate states are far from the boundary.

The corresponding Euler negativity probabilities are strictly positive, because a Gaussian increment has full support, and they should not be called zero. They are, however, extraordinarily small. \Cref{tab:cir_negativity} reports them at $h=1/520$, corresponding to weekly data with $M=10$. At the long-run means the standardised distance to the boundary ranges from $111$ to $203$ standard deviations, so the one-step probabilities are of order $10^{-2667}$ to $10^{-8986}$. Even from the implausibly low state $r=0.001$ the smallest distance is $13.8$ standard deviations and the largest probability is about $2\times10^{-43}$. All but the last of these underflow to exactly zero in double precision, which is why the table reports base-ten logarithms rather than the probabilities themselves.

The assumptions of the generic theorem are genuinely violated by \eqref{eq:cir_rates}, and the estimator applied is therefore not the estimator analysed. But for $r$ and $r^*$ that gap is a matter of proof coverage, not of numerical damage: a violation probability of $10^{-2667}$ per step is not a phenomenon any simulation will encounter.

\begin{table}[t]
\centering
\caption{One-step Euler negativity diagnostics for the square-root interest-rate processes at $h=1/520$, using the estimates of Table 3 of \citet{BrandtSantaClara2002}. The probability is strictly positive at every state; ``underflow'' records whether direct double-precision evaluation returns exactly zero. Base-ten logarithms are reported because the probabilities are not representable.}
\label{tab:cir_negativity}
\small
\begin{tabular}{llrrrr}
\toprule
Process & State & $r$ & $2\lambda\vartheta/\sigma^2$ & $z$ & $\log_{10}\Pp(\text{negative})$ \\
\midrule
U.S. (vs. U.K.) & long-run mean & 0.053 & 38.4 & 187.5 & $-7636$ \\
 & half the mean & 0.0265 & & 132.6 & $-3823$ \\
 & 1 per cent & 0.01 & & 81.6 & $-1449$ \\
 & 0.1 per cent & 0.001 & & 26.5 & $-154$ \\
\midrule
U.K. & long-run mean & 0.074 & 22.9 & 110.8 & $-2667$ \\
 & half the mean & 0.037 & & 78.4 & $-1337$ \\
 & 1 per cent & 0.01 & & 41.0 & $-366$ \\
 & 0.1 per cent & 0.001 & & 13.8 & $-42.6$ \\
\midrule
U.S. (vs. Germany) & long-run mean & 0.058 & 48.5 & 203.4 & $-8986$ \\
 & 0.1 per cent & 0.001 & & 27.6 & $-167$ \\
\midrule
Germany & long-run mean & 0.064 & 6.4 & 137.4 & $-4099$ \\
 & 0.1 per cent & 0.001 & & 17.4 & $-67.0$ \\
\bottomrule
\end{tabular}
\end{table}

The market-incompleteness state behaves differently. Write $H_t=\epsilon_t^2$ for it; the letter $X$ is reserved for the Brownian motion driving the log exchange rate in \eqref{eq:corr_block}. Equation (37) of \citet{BrandtSantaClara2002} specifies
\begin{equation}\label{eq:epsilon_cir}
dH_t=(\beta^2-2\alpha H_t)dt+2\beta\sqrt{H_t}\,dY_t,
\end{equation}
which is a CIR process with $\kappa=2\alpha$, long-run mean $\beta^2/(2\alpha)$, and diffusion coefficient $2\beta$. Its Feller ratio is
\[
\frac{2\kappa\vartheta_X}{\sigma_X^2}
=\frac{2\beta^2}{4\beta^2}
=\frac12
\]
for every $\alpha>0$ and every $\beta\ne0$. The ratio is exactly one half independently of the parameter values, so no estimate can move this process away from its boundary: zero is attainable, and the instantaneous covariance degenerates there. This is a structural feature of the specification, not a property of a particular estimate. \Cref{prop:min_incomplete} below shows that the boundary is not merely attainable in principle but attained exactly in the reported results.

\begin{table}[t]
\centering
\caption{Feller ratios at the estimates reported in Table 3 of \citet{BrandtSantaClara2002}. For a process $dZ=\kappa(\vartheta_Z-Z)dt+\sigma_Z\sqrt Z\,dB$ the ratio is $2\kappa\vartheta_Z/\sigma_Z^2$, and the origin is unattainable if and only if it is at least one. For the incompleteness state the ratio equals $1/2$ for all admissible parameters.}
\label{tab:feller}
\begin{tabular}{llrrrr}
\toprule
System & State & $\kappa$ & $\vartheta_Z$ & $\sigma_Z$ & $2\kappa\vartheta_Z/\sigma_Z^2$ \\
\midrule
U.S. vs. U.K. & $r$ (U.S.) & 0.284 & 0.053 & 0.028 & 38.4 \\
              & $r^*$ (U.K.) & 0.486 & 0.074 & 0.056 & 22.9 \\
              & $H=\epsilon^2$ & 0.640 & 0.0121 & 0.176 & 0.5 \\
U.S. vs. Germany & $r$ (U.S.) & 0.305 & 0.058 & 0.027 & 48.5 \\
              & $r^*$ (Germany) & 0.088 & 0.064 & 0.042 & 6.4 \\
              & $H=\epsilon^2$ & 0.676 & 0.0151 & 0.202 & 0.5 \\
\bottomrule
\end{tabular}
\end{table}

\section{The simulated state and the reality of the incompleteness process}\label{app:reality}

The previous subsection treats \eqref{eq:epsilon_cir} as a state in its own right, which is how the model presents it. The estimator does not. Equation (31) and Appendix B of \citet{BrandtSantaClara2002} eliminate $\epsilon_t$ using the variance identity and simulate the four-dimensional state
\begin{equation}\label{eq:simulated_state}
(r_t,\;r_t^*,\;e_t,\;v_t),
\end{equation}
where $e_t$ is the log exchange rate and $v_t$ is its instantaneous volatility, proxied by a one-week at-the-money implied volatility. The incompleteness state is therefore never discretised directly. This relocates the boundary problem rather than removing it, and the relocated version is the one that binds.

The identity that eliminates $\epsilon_t$ is
\begin{equation}\label{eq:epsilon_recovered}
\epsilon_t
=\sqrt{
v_t^2
-(\phi_t^2+\phi_t^{*2}-2\rho_{ww^*}\phi_t\phi_t^*)
-(\psi_t^2+\psi_t^{*2}-2\rho\,\psi_t\psi_t^*)
},
\end{equation}
where $\phi_t=\phi\sqrt{r_t}$ and $\phi_t^*=\phi^*\sqrt{r_t^*}$ are the market prices of domestic and foreign interest-rate risk and $\psi_t,\psi_t^*$ are the market prices of pure currency risk. The state $\epsilon_t$ appears in the diffusion of the log exchange rate, through $v_t\,dX_t=\phi_t dW_t-\phi_t^*dW_t^*+\psi_t dZ_t-\psi_t^*dZ_t^*+\epsilon_t dU_t$, and in the coefficients of the $v_t$ equation obtained from It\^o's lemma. Every Euler substep of \eqref{eq:simulated_state} therefore requires the radicand of \eqref{eq:epsilon_recovered} to be nonnegative at the simulated state, not merely at the observed data. Ordinary Euler discretisation of a system whose coefficients are defined by \eqref{eq:epsilon_recovered} offers no such guarantee: the increments of $v_t$ are conditionally Gaussian and unbounded below, while the subtracted terms are continuous and generically bounded away from zero.

This is the operative form of the degeneracy, and three features make it sharp rather than hypothetical.

First, the constraint is not generically slack. Under the conditions of \Cref{prop:min_incomplete}, in the case where a nonnegativity constraint is the active one, the identification rule makes $\min_t\epsilon_t^2$ equal to zero exactly, so there is at least one sample date at which the reconstructed incompleteness state sits on the boundary. The transition simulated \emph{out of} that date therefore starts on the boundary, unless the active date happens to be the terminal observation, in which case no transition is simulated from it. This is a conclusion from the optimisation geometry, not from the reported table: the minima of $0.0000$ in Table 4 of \citet{BrandtSantaClara2002} are rounded to four decimal places and are corroboration rather than proof. Nothing here says that every simulated path starts from the boundary, only that at least one does.

Second, the failure is not a recoverable numerical event such as a negative variance that can be floored at zero. It makes the drift and diffusion coefficients of the next substep complex-valued, so the chain is undefined rather than inaccurate.

Third, the paper reports assigning zero likelihood to parameter values that violate $\epsilon_t^2\ge0$ at the observed dates. That is a restriction on parameters, not a rule for what the simulator does when an \emph{intermediate} Euler state violates it, and the two are different objects: the constraint is checked at $N$ observation dates, whereas the chain visits $NMS$ intermediate states.

As with the interest rates, a truncation, reflection, absorption, or exact-transition rule would resolve the issue. Each such rule defines a \emph{different} approximating chain, with a different transition density and therefore a different estimator, and each requires its own convergence analysis; they are not interchangeable implementation details. The published paper does not specify one.

The size of the problem can be computed exactly on the incompleteness specification \eqref{eq:epsilon_cir} taken at face value. For an Euler step of size $h$ the probability of a negative next state from $H_t=x>0$ is
\begin{equation}\label{eq:negative_prob}
P_h(x)=\Phi\!\left(
-\frac{x+(\beta^2-2\alpha x)h}{2\beta\sqrt{xh}}
\right).
\end{equation}
This vanishes at $x=0$, where the diffusion term vanishes while the drift $\beta^2h$ is strictly positive, and again for large $x$, so it has an interior maximum. That maximum has a closed form, and it is not what one would hope.

\begin{proposition}[Exact worst case of the one-step violation probability]\label{prop:worst_case_boundary}
Let $\alpha>0$, $\beta>0$ and $0<h<1/(2\alpha)$. Then $P_h$ attains its supremum on $(0,\infty)$ at
\begin{equation}\label{eq:worst_case_argmax}
x_h^*=\frac{\beta^2h}{1-2\alpha h},
\end{equation}
and the supremum equals
\begin{equation}\label{eq:worst_case}
\sup_{x>0}P_h(x)=\Phi\!\left(-\sqrt{1-2\alpha h}\right).
\end{equation}
Consequently $x_h^*\sim\beta^2h$ and $\sup_xP_h(x)\to\Phi(-1)\approx0.1587$ as $h\downarrow0$, and the limiting worst case depends on neither $\alpha$ nor $\beta$.
\end{proposition}

The sign convention matters here. We take $\beta>0$, which matches the reported estimates; for $\beta<0$ the process is unchanged in law and every occurrence of $\beta$ in the standard deviation should be read as $|\beta|$, after which the statement is identical with $\beta^2=|\beta|^2$ throughout.

\begin{proof}
Write $a=1-2\alpha h>0$ and $b=\beta^2h>0$. The argument of $\Phi$ in \eqref{eq:negative_prob} is $-g(x)$ with
\[
g(x)=\frac{x+(\beta^2-2\alpha x)h}{2\beta\sqrt{xh}}=\frac{ax+b}{2\beta\sqrt h\,\sqrt x},
\]
so, $\Phi$ being strictly increasing, maximising $P_h$ is minimising $g$ over $x>0$. Differentiating the numerator of $g$ with respect to $\sqrt x$, or directly,
\[
g'(x)=\frac{1}{2\beta\sqrt h}\cdot\frac{a\sqrt x-(ax+b)/(2\sqrt x)}{x}
=\frac{ax-b}{4\beta\sqrt h\,x^{3/2}},
\]
which is negative for $x<b/a$, zero at $x=b/a$ and positive for $x>b/a$. Hence $g$ has a unique minimum at $x_h^*=b/a$, which is \eqref{eq:worst_case_argmax}. There
\[
g(x_h^*)=\frac{a(b/a)+b}{2\beta\sqrt h\,\sqrt{b/a}}=\frac{2b}{2\beta\sqrt h\,\sqrt{b/a}}=\frac{\sqrt{ab}}{\beta\sqrt h}=\frac{\beta\sqrt{ah}}{\beta\sqrt h}=\sqrt a,
\]
giving \eqref{eq:worst_case}. The limits follow since $a\to1$ as $h\downarrow0$.
\end{proof}

The practical reading is that halving the Euler step does not make the worst case safer. It moves the dangerous region towards zero, in proportion to $h$, while the peak violation probability rises very slightly towards $\Phi(-1)$ from below. \Cref{tab:boundary_max} reports \eqref{eq:worst_case_argmax} and \eqref{eq:worst_case} at the reported estimates.

Two qualifications keep this in proportion. First, $\sup_xP_h(x)$ is a worst case over states, not a property of a typical path: it says that \emph{if} the chain is at $x_h^*$ then the next step is negative with probability about $0.16$, and says nothing about how often the chain is there. Second, the dangerous region shrinks with $h$, so a finer discretisation does reduce the time spent in it. Whether the overall risk falls therefore depends on the invariant law of the chain near the boundary, which the discretisation alone does not determine. What \Cref{prop:worst_case_boundary} rules out is the comfortable assumption that refinement makes the per-step boundary problem itself disappear.

\Cref{fig:negative_euler} plots \eqref{eq:negative_prob} at these step sizes. It is an exact illustration of the class of failure, not a simulation of the implemented chain, whose boundary is \eqref{eq:epsilon_recovered} and whose exact crossing probability depends on the drift and diffusion of $v_t$ derived in Appendix B of \citet{BrandtSantaClara2002}.

\begin{table}[t]
\centering
\caption{Worst-case one-step Euler negativity probability for the incompleteness specification \eqref{eq:epsilon_cir}, at the estimates of Table 3 of \citet{BrandtSantaClara2002}. The step $h=1/520$ corresponds to weekly data with $M=10$ and $h=1/1040$ to $M=20$; the third row is a coarser illustration. The maximum is essentially invariant to $h$, in accordance with \Cref{prop:worst_case_boundary}, while the maximising state scales with $h$.}
\label{tab:boundary_max}
\begin{tabular}{llrrr}
\toprule
System & Step & $H^*_h$ & $\epsilon^*_h=\sqrt{H^*_h}$ & $\max_xP_h(x)$ \\
\midrule
U.S. vs. U.K.
 & $h=1/520$ & $1.49\times10^{-5}$ & 0.00386 & 0.1588 \\
 & $h=1/1040$ & $7.45\times10^{-6}$ & 0.00273 & 0.1587 \\
 & $h=0.02$ & $1.57\times10^{-4}$ & 0.01253 & 0.1602 \\
\midrule
U.S. vs. Germany
 & $h=1/520$ & $1.96\times10^{-5}$ & 0.00443 & 0.1588 \\
 & $h=1/1040$ & $9.82\times10^{-6}$ & 0.00313 & 0.1587 \\
 & $h=0.02$ & $2.07\times10^{-4}$ & 0.01438 & 0.1603 \\
\bottomrule
\end{tabular}
\end{table}

\section{Invertibility of the volatility--incompleteness map}\label{app:invert}

The model can be written with either volatility or incompleteness as the primitive state, and the two representations are equivalent only if the map between them is well defined in both directions. This subsection analyses that map. It is not a claim about what the implemented estimator computes: as described in \Cref{app:reality}, the implementation takes $v_t$ as the primitive simulated state and uses the variance identity in the easy direction, to reconstruct $\epsilon_t$. No quadratic is solved at any Euler step.

The time-varying pure-currency risk prices have the form
\begin{equation}\label{eq:risk_price_v}
\psi_t=a_t+bv_t,
\qquad
\psi_t^*=a_t^*+b^*v_t,
\end{equation}
where $a_t$ and $a_t^*$ depend on the interest-rate differential and the log exchange rate. The variance identity is
\begin{equation}\label{eq:vol_identity}
v_t^2
=D_t+\psi_t^2+(\psi_t^*)^2-2\rho\psi_t\psi_t^*,
\end{equation}
where $D_t$ collects the interest-rate risk term and $\epsilon_t^2$, and $\rho=\rho_{zz^*}$.

Substituting \eqref{eq:risk_price_v} into \eqref{eq:vol_identity} and collecting terms defines
\begin{equation}\label{eq:F_definition}
F(v;\,\text{state})=(1-\kappa_b)v^2-2\eta_tv-\zeta_t,
\end{equation}
so that the identity reads $F(v_t;\text{state})=0$, where
\begin{align}
\kappa_b&=b^2+(b^*)^2-2\rho bb^*, \label{eq:kappa_b}\\
\eta_t&=a_tb+a_t^*b^*-\rho(a_tb^*+a_t^*b),\\
\zeta_t&=D_t+a_t^2+(a_t^*)^2-2\rho a_ta_t^*.
\end{align}
Both $\kappa_b$ and the state-dependent part of $\zeta_t$ are quadratic forms in the two-dimensional correlation matrix $\bigl(\begin{smallmatrix}1&\rho\\\rho&1\end{smallmatrix}\bigr)$, evaluated at $(b,-b^*)$ and $(a_t,-a_t^*)$ respectively. When that matrix is positive semidefinite and $D_t\ge0$ one therefore has $\kappa_b\ge0$ and $\zeta_t\ge0$.

The two directions of the map are not symmetric.

\emph{Volatility primitive.} Given $v_t$ and the remaining states, \eqref{eq:vol_identity} determines $\epsilon_t^2$ by subtraction, uniquely and without any branch condition. This is the direction the estimator uses, and it is unproblematic apart from the reality requirement discussed in \Cref{app:reality}.

\emph{Incompleteness primitive.} Given $\epsilon_t^2$ and the remaining states, recovering a positive $v_t$ requires solving $F(v;\text{state})=0$. This is the direction implied by reading equations (30) and (37) of \citet{BrandtSantaClara2002} as the primitive specification, with the volatility dynamics of their equation (31) derived from it by It\^o's lemma. The implicit function theorem gives local solvability wherever
\begin{equation}\label{eq:implicit_derivative}
\frac{\partial F}{\partial v}(v;\text{state})=2(1-\kappa_b)v-2\eta_t\ne0,
\end{equation}
that is away from the singular set
\begin{equation}\label{eq:singular_set}
v=\frac{\eta_t}{1-\kappa_b}
\qquad(\kappa_b\ne1),
\end{equation}
at which the two roots coalesce and the local inverse fails to exist. Local solvability is not global uniqueness, and the following proposition describes the global branch structure.

The following proposition concerns the state transformation between volatility and incompleteness. It is not a statement about the implemented optimisation algorithm, which uses the identity only in the direction described above.

\begin{proposition}[Branch structure of the volatility identity]\label{prop:vol_branch}
Let $\zeta_t>0$.
\begin{enumerate}[label=(\roman*)]
\item If $\kappa_b<1$, equation $F(v;\text{state})=0$ has exactly one positive root and one negative root, so the positive branch is unique and the inverse map is globally single valued on the positive half-line.
\item If $\kappa_b>1$, the equation has either no positive root or exactly two positive roots. A unique positive root occurs only in the degenerate case $\eta_t^2=(\kappa_b-1)\zeta_t$ of a repeated root, which is exactly the singular set \eqref{eq:singular_set}. Positive roots exist if and only if $\eta_t<0$ and $\eta_t^2\ge(\kappa_b-1)\zeta_t$, and are then
\[
v_t^{\pm}=\frac{-\eta_t\pm\sqrt{\eta_t^2-(\kappa_b-1)\zeta_t}}{\kappa_b-1}.
\]
\end{enumerate}
\end{proposition}

\begin{proof}
Rewrite $F(v;\text{state})=0$ as $(\kappa_b-1)v^2+2\eta_tv+\zeta_t=0$.

(i) If $1-\kappa_b>0$ the product of the roots is $-\zeta_t/(1-\kappa_b)<0$, so the real roots have opposite signs. The discriminant $4\{\eta_t^2+(1-\kappa_b)\zeta_t\}$ is strictly positive, so both roots are real and exactly one is positive.

(ii) If $\kappa_b-1>0$ the product of the roots is $\zeta_t/(\kappa_b-1)>0$ and their sum is $-2\eta_t/(\kappa_b-1)$. A positive product excludes roots of opposite signs and excludes a zero root, so real roots are either both positive or both negative, and both are positive exactly when the sum is positive, that is when $\eta_t<0$. Real roots exist when the discriminant $4\{\eta_t^2-(\kappa_b-1)\zeta_t\}$ is nonnegative, and the stated formula is the quadratic formula. The roots coincide when $\eta_t^2=(\kappa_b-1)\zeta_t$, in which case the common value is $-\eta_t/(\kappa_b-1)=\eta_t/(1-\kappa_b)$, which is \eqref{eq:singular_set}.
\end{proof}

\Cref{tab:implicit_coefficients} evaluates $\kappa_b$ at the model-C estimates of Table 4 of \citet{BrandtSantaClara2002}. The two systems fall on opposite sides of the threshold.

For the U.S.--Germany system, $\kappa_b\approx0.0702<1$, case (i) applies, and the inverse map is globally single valued: given any admissible incompleteness state there is exactly one positive volatility consistent with it.

For the U.S.--U.K. system, $\kappa_b\approx1.9982>1$ and case (ii) applies. The reported coefficients therefore do not define a globally unique inverse mapping from the incompleteness state to positive volatility without a branch-selection condition. At states where positive roots exist there are two of them, and the specification does not say which is intended; at states where $\eta_t\ge0$ or $\eta_t^2<(\kappa_b-1)\zeta_t$ there is none. Because $\zeta_t\ge0$ follows from positive semidefiniteness of the two-dimensional correlation matrix together with $D_t\ge0$, both of which hold by construction, this is a consequence of $\kappa_b>1$ alone and not of a knife-edge parameter value.

Two qualifications matter for how much weight this carries. It does not affect the reconstruction of $\epsilon_t$ from an observed implied volatility, which is the direction the estimation uses and which is unique. And it does not by itself show that the reported estimates are wrong. What it shows is that the two representations of the model are not interchangeable at the U.S.--U.K. estimates, so that reading the specification as a generative diffusion in the incompleteness state requires an additional branch-selection rule that the paper does not supply. Appendix B of \citet{BrandtSantaClara2002} defines the drift and diffusion of $v_t$ only implicitly, with those coefficients appearing inside the formulas used to define them, so a well-posedness condition is required there in any case.

\begin{table}[t]
\centering
\caption{Quadratic coefficient implied by the model-C volatility loadings reported in Table 4 of \citet{BrandtSantaClara2002}, where $\kappa_b=b^2+(b^*)^2-2\rho bb^*$ with $b=c_3$, $b^*=c_3^*$ and $\rho=\rho_{zz^*}$. The sign of $1-\kappa_b$ determines which case of \Cref{prop:vol_branch} applies.}
\label{tab:implicit_coefficients}
\begin{tabular}{lrrrrrl}
\toprule
Market & $b$ & $b^*$ & $\rho$ & $\kappa_b$ & $1-\kappa_b$ & Branch structure \\
\midrule
U.S. vs. U.K. & 1.514 & 2.581 & 0.89 & 1.9982 & $-0.9982$ & zero or two positive roots \\
U.S. vs. Germany & $-0.142$ & 0.127 & 0.94 & 0.0702 & 0.9298 & unique positive root \\
\bottomrule
\end{tabular}
\end{table}

\section{Feasibility of the Brownian correlation matrix}\label{app:psd}

The application specifies an instantaneous correlation matrix for the four Brownian motions $(W_t,W_t^*,X_t,Y_t)$, where $W$ and $W^*$ drive the two interest rates, $X$ drives the log exchange rate and $Y$ drives the incompleteness state. Write it in block form as
\begin{equation}\label{eq:corr_block}
R_t=
\begin{pmatrix}
A_t&r_t\\
r_t^{\mathsf T}&1
\end{pmatrix},
\end{equation}
where $A_t$ is the $3\times3$ block for $(W_t,W_t^*,X_t)$ and $r_t=(\rho_{wy},\rho_{w^*y},\rho_{xy})^{\mathsf T}$ collects the three correlations with $Y_t$. Three entries of $R_t$ are free parameters, namely the components of $r_t$; the remaining off-diagonal entries are implied. From equation (28) of \citet{BrandtSantaClara2002},
\begin{equation}\label{eq:implied_corr}
\rho_{wx,t}=\frac{\phi_t-\rho_{ww^*}\phi_t^*}{v_t},
\qquad
\rho_{w^*x,t}=\frac{\rho_{ww^*}\phi_t-\phi_t^*}{v_t},
\end{equation}
state dependent through $\phi_t=\phi\sqrt{r_t}$, $\phi_t^*=\phi^*\sqrt{r_t^*}$ and $v_t$.

Three different questions must be kept apart, and the three subsections below answer them separately. First, is the matrix valid at the reported estimates and at reference states the calibrated model can reach? Second, does the specification guarantee validity for every admissible parameter and state? Third, did the numerical implementation impose safeguards not described in the published paper? The answers are yes, no, and unknown.

\subsection{The feasible state space}\label{app:psd_feasible}

Everything here depends on which states the model can reach, so we begin with the complete variance identity. Equation (29) of \citet{BrandtSantaClara2002} is
\begin{equation}\label{eq:full_variance_identity}
v_t^2=Q_t+C_t+\epsilon_t^2,
\end{equation}
where
\[
Q_t=\phi_t^2+\phi_t^{*2}-2\rho_{ww^*}\phi_t\phi_t^*,
\qquad
C_t=\psi_t^2+\psi_t^{*2}-2\rho_{zz^*}\psi_t\psi_t^*,
\]
are the interest-rate and currency-risk quadratic forms and $\epsilon_t^2\ge0$ is the incompleteness state. Both quadratic forms are nonnegative whenever the corresponding two-dimensional correlation matrix is positive semidefinite. A state is feasible precisely when \eqref{eq:full_variance_identity} holds with $\epsilon_t^2\ge0$, that is when
\begin{equation}\label{eq:true_floor}
v_t^2\ \ge\ Q_t+C_t .
\end{equation}
The weaker bound $v_t^2\ge Q_t$ is necessary but not sufficient, and it is attained only in the degenerate configuration $C_t=\epsilon_t^2=0$.

For the constant-risk-price specification, $C_t$ is a constant that can be read directly from the published estimates. Table 3 of \citet{BrandtSantaClara2002} reports $\psi^2-\rho_{zz^*}\psi\psi^*$ and $\psi^{*2}-\rho_{zz^*}\psi\psi^*$ separately, and their sum is exactly $C$. This gives
\[
C_{\text{U.K.}}=0.024+(-0.021)=0.003,
\qquad
C_{\text{Germany}}=(-0.010)+0.013=0.003 .
\]
Both are strictly positive, so at the reported estimates the degenerate configuration is \emph{not} reachable: the feasible minimum volatility exceeds $\sqrt{Q_t}$ by a margin that does not vanish. At the long-run mean rates the feasible minimum is $0.0640$ for the U.S.--U.K. system against $\sqrt{Q_t}=0.0330$, and $0.0622$ against $0.0295$ for the U.S.--Germany system.

For the time-varying specifications $C_t$ is state dependent, and in model C it depends on $v_t$ itself through \eqref{eq:risk_price_v}. That looks like an obstruction and is not one, for two reasons.

First, model B carries no volatility loading: Table 4 reports three coefficients per risk price for that specification against four for model C, so under model B the quantity $C_t$ is an explicit function of the interest-rate differential and the log exchange rate, and $v_t$ follows from \eqref{eq:full_variance_identity} by taking a square root.

Second, the apparent circularity in model C is exactly the quadratic of \Cref{app:invert}. Substituting \eqref{eq:risk_price_v} into $C_t$ and then into \eqref{eq:full_variance_identity} gives
\begin{equation}\label{eq:model_c_fixed_point}
(1-\kappa_b)v_t^2-2\eta_tv_t-\left(Q_t+\zeta^a_t+\epsilon_t^2\right)=0,
\qquad
\zeta^a_t=a_t^2+a_t^{*2}-2\rho a_ta_t^*,
\end{equation}
so $v_t$ is \emph{determined} by $(r_t,r_t^*,e_t,\epsilon_t^2)$ rather than being an independent unknown. The fixed point is solvable in closed form, and by \Cref{prop:vol_branch} it has one positive solution when $\kappa_b<1$ and either none or two when $\kappa_b>1$.

What remains genuinely unobserved is only the level of the log exchange rate $e_t$. Rather than requiring its realised path we sweep it, which yields a statement about the reachable state space rather than about one realisation, and is therefore stronger than a check against the observed series would be. \Cref{tab:time_varying} reports the result over a grid of $56{,}784$ states per specification, spanning interest rates from $0.05$ to $5$ times their long-run means, $e_t\in[-2,2]$, and $\epsilon_t^2\in\{0,0.002,0.01,0.05\}$. Because Table 4 does not re-report the interest-rate risk prices or the three free correlations, those are taken from Table 3; that is an assumption, and it is the only one.

Three of the four specifications are unproblematic. For model B in both systems and for model C in the U.S.--Germany system, every swept state admits exactly one positive volatility and the correlation matrix is positive definite at all of them, with worst minimum eigenvalues of $0.070$, $0.135$ and $0.088$ respectively.

The U.S.--U.K. model C specification behaves differently, and the behaviour sharpens \Cref{prop:vol_branch} from a statement about branch counts into a statement about how much of the state space is covered at all. Of the $56{,}784$ states swept, $37{,}536$, or $66\%$, admit \emph{no positive volatility whatsoever}; the remaining $19{,}248$ admit exactly two. Not one state admits exactly one, which is \Cref{prop:vol_branch}(ii) confirmed as an empirical fact rather than a possibility. Where a branch does exist the correlation matrix is comfortably valid, with minimum eigenvalue at least $0.445$, so this is a well-posedness failure rather than a correlation failure. Moreover the region shrinks as incompleteness grows: at the long-run mean rates a branch exists only for $e_t\gtrsim0.3$ when $\epsilon_t^2=0$, and for no $e_t$ in the swept range once $\epsilon_t^2=0.05$. Read as a generative diffusion in the incompleteness state, that specification is therefore undefined on most of the state space it would need to traverse.

\begin{table}[t]
\centering
\caption{Feasibility of the time-varying specifications over a swept state space, using the Table 4 loadings together with the Table 3 interest-rate risk prices and correlations. Volatility is obtained from \eqref{eq:full_variance_identity}, through \eqref{eq:model_c_fixed_point} under model C, so every matrix evaluated satisfies the variance identity by construction. No state with an indefinite matrix was found in any specification.}
\label{tab:time_varying}
\small
\begin{tabular}{llrrrrr}
\toprule
System & Model & States & No branch & Two branches & Matrices & $\min\lambda_{\min}$ \\
\midrule
U.S. vs. U.K. & B & 56{,}784 & 0 & 0 & 56{,}784 & $+0.0699$ \\
 & C & 56{,}784 & 37{,}536 & 19{,}248 & 38{,}496 & $+0.4455$ \\
\midrule
U.S. vs. Germany & B & 56{,}784 & 0 & 0 & 56{,}784 & $+0.1351$ \\
 & C & 56{,}784 & 0 & 0 & 56{,}784 & $+0.0880$ \\
\bottomrule
\end{tabular}
\end{table}

\subsection{Validity at reported reference states}\label{app:psd_reported}

\Cref{tab:correlation_audit} evaluates $R_t$ across reference states, reporting the complete decomposition \eqref{eq:full_variance_identity} so that feasibility can be judged from the identity rather than from the insufficient bound.

At every feasible state examined the matrix is positive definite. At the long-run means with the sample mean volatility the minimum eigenvalue is $0.683$ for the U.S.--U.K. system and $0.639$ for the U.S.--Germany system, with Schur complements of $0.997$ and $0.994$. Pushing to the feasible minimum volatility at rates ten times their long-run means still leaves minimum eigenvalues of $0.115$ and $0.125$. Perturbing the four estimated correlations independently by their reported standard errors, in $5{,}000$ draws per system, produced no violation of positive semidefiniteness and no entrywise violation.

A systematic search reinforces this. Sweeping both interest rates over multiples of their long-run means from zero to one hundred, and the incompleteness state over the values $\epsilon_t^2\in\{0,0.005,0.02,0.10\}$, always taking $v_t$ from the identity \eqref{eq:full_variance_identity}, the smallest minimum eigenvalue encountered is $+0.0128$ for the U.S.--U.K. system and $+0.0130$ for the U.S.--Germany system, both at the extreme corner of the sweep where the interest rates are a hundred times their long-run means. \textbf{No model-feasible state with an indefinite correlation matrix was found.} We therefore make no claim that the calibrated model reaches an invalid correlation matrix.

The last row of each block in \Cref{tab:correlation_audit} records what happens at $v_t=\sqrt{Q_t}$, which an earlier version of this analysis mistakenly treated as the floor. There the matrix is indeed indefinite, with minimum eigenvalues $-0.0009$ and $-0.0027$, while every entry remains inside $[-1,1]$. But the implied $\epsilon_t^2$ equals $-C=-0.003$, so the state violates \eqref{eq:full_variance_identity} and lies outside the model. It is reported here only to mark the distinction, not as evidence.

\subsection{Global validity of the specification}\label{app:psd_global}

That the reported estimates are valid is not the same as the specification guaranteeing validity. Two results bear on this, and both are conditional algebraic statements rather than claims about the calibrated model.

\begin{proposition}[Entrywise validity is automatic]\label{prop:entrywise}
Whenever $v_t^2\ge Q_t$, which \eqref{eq:full_variance_identity} always implies, $|\rho_{wx,t}|\le1$ and $|\rho_{w^*x,t}|\le1$.
\end{proposition}

\begin{proof}
$Q_t-(\phi_t-\rho_{ww^*}\phi_t^*)^2=\phi_t^{*2}(1-\rho_{ww^*}^2)\ge0$, so $(\phi_t-\rho_{ww^*}\phi_t^*)^2\le Q_t\le v_t^2$, giving $|\rho_{wx,t}|\le1$. The second bound is symmetric, using $Q_t-(\rho_{ww^*}\phi_t-\phi_t^*)^2=\phi_t^2(1-\rho_{ww^*}^2)\ge0$.
\end{proof}

The consequence is that the entrywise restriction the paper imposes, through a logistic reparameterisation of the correlations, is not doing any work on these three entries: the model structure already implies it. What it does not imply is the joint condition. When $A_t\succ0$ that condition is the Schur complement
\begin{equation}\label{eq:schur}
1-r_t^{\mathsf T}A_t^{-1}r_t\ge0,
\end{equation}
and constraining each entry of $r_t$ to $[-1,1]$ is not sufficient for it.

The remaining question is what happens when $A_t$ is singular, where \eqref{eq:schur} is unavailable because $A_t^{-1}$ does not exist. This occurs exactly in the degenerate configuration identified above.

\begin{proposition}[Compatibility in the degenerate configuration]\label{prop:floor_overdetermined}
Suppose $C_t=\epsilon_t^2=0$ and $Q_t>0$, so that $v_t^2=Q_t$. Then $A_t$ is singular with a one-dimensional kernel, $X_t$ is a deterministic linear combination of $W_t$ and $W_t^*$, and $R_t\succeq0$ holds if and only if
\begin{equation}\label{eq:induced_rho_xy}
\rho_{xy}=\frac{\phi_t\rho_{wy}-\phi_t^*\rho_{w^*y}}{v_t} .
\end{equation}
Any other value makes $R_t$ indefinite, however far each individual entry lies inside $[-1,1]$.
\end{proposition}

\begin{proof}
If $C_t=\epsilon_t^2=0$ then equation (28) reduces to $v_t\,dX_t=\phi_t\,dW_t-\phi_t^*\,dW_t^*$, so $X_t$ lies in the span of $(W_t,W_t^*)$ and $A_t$ is singular. Let $u=(\phi_t,-\phi_t^*,-v_t)^{\mathsf T}$, which spans the kernel of $A_t$, so $A_tu=0$ and $u^{\mathsf T}A_tu=0$. Extend it to $\tilde u=(u^{\mathsf T},0)^{\mathsf T}\in\R^4$ and let $e_4$ be the fourth coordinate vector. For $t\in\R$,
\[
(\tilde u+te_4)^{\mathsf T}R_t(\tilde u+te_4)
=u^{\mathsf T}A_tu+2t\,u^{\mathsf T}r_t+t^2
=2t\,u^{\mathsf T}r_t+t^2 .
\]
If $u^{\mathsf T}r_t\ne0$, choosing $t$ of small magnitude and of sign opposite to $u^{\mathsf T}r_t$ makes this expression negative, so $R_t$ is not positive semidefinite. Conversely if $u^{\mathsf T}r_t=0$ the quadratic form is $t^2\ge0$ along that plane, and positive semidefiniteness of $R_t$ follows from that of $A_t$ together with the Moore--Penrose compatibility condition $1-r_t^{\mathsf T}A_t^{+}r_t\ge0$, which holds because $r_t$ then lies in the range of $A_t$. Writing out $u^{\mathsf T}r_t=\phi_t\rho_{wy}-\phi_t^*\rho_{w^*y}-v_t\rho_{xy}=0$ gives \eqref{eq:induced_rho_xy}.
\end{proof}

This says that in the degenerate configuration the correlation $\rho_{xy}$ is not a free parameter but a function of the others, whereas the application estimates it freely. It is a genuine gap in the specification: nothing in the model enforces \eqref{eq:induced_rho_xy}, and no Cholesky, partial-correlation, hyperspherical or Schur-complement parameterisation is stated that would enforce \eqref{eq:schur} more generally.

It is not, however, a defect in the reported estimates, for the reason established in \Cref{app:psd_feasible}: the configuration requires $C_t=0$, and at the reported estimates $C=0.003>0$. To see how far the reported values are from satisfying \eqref{eq:induced_rho_xy}, evaluating the right-hand side at three times the long-run rates gives $-0.0538$ and $-0.0696$ against reported $\rho_{xy}$ of $-0.012$ and $+0.006$; but that comparison is made at a state the calibrated model does not reach, and it establishes only that the identity would have to be imposed rather than estimated if the specification were used in a regime where $C_t$ approached zero.

\subsection{Unknown implementation safeguards}\label{app:psd_implementation}

Whether the numerical implementation rejected invalid parameter draws, projected them onto the positive-semidefinite cone, worked with a Cholesky factor, or took some other precaution, is not described in the published paper and cannot be determined from it. We record this as unknown. It is entirely possible that the optimiser never encountered an invalid matrix, either because the safeguards existed or because, as \Cref{app:psd_reported} suggests, the feasible region around the estimates is comfortably interior. Nothing in this appendix should be read as a claim about what the implementation did.

\begin{table}[t]
\centering
\caption{Feasibility and validity of the four-dimensional Brownian correlation matrix under the complete variance identity $v_t^2=Q_t+C_t+\epsilon_t^2$, at the constant-risk-price estimates of Table 3 of \citet{BrandtSantaClara2002}. Rates are multiples of their long-run means. A state is feasible when the implied $\epsilon_t^2$ is nonnegative. The final row of each block uses $v_t=\sqrt{Q_t}$, which requires $C_t=0$ and is therefore infeasible at these estimates; it is shown to separate what the model excludes from what it permits.}
\label{tab:correlation_audit}
\footnotesize
\setlength{\tabcolsep}{3pt}
\begin{tabular}{llrrrrcrr}
\toprule
System & State & $Q_t$ & $C_t$ & $\epsilon_t^2$ & $v_t$ & Feas. & $\max|\rho_{ij}|$ & $\lambda_{\min}(R_t)$ \\
\midrule
U.S. vs. U.K.
 & means, $v_t=\bar v$ & 0.0011 & 0.0030 & $0.0065$ & 0.1030 & yes & 0.313 & $+0.6834$ \\
 & means, $v_t$ min & 0.0011 & 0.0030 & $0.0000$ & 0.0640 & yes & 0.503 & $+0.4875$ \\
 & rates $\times2$, at min & 0.0022 & 0.0030 & $0.0000$ & 0.0720 & yes & 0.633 & $+0.3540$ \\
 & rates $\times3$, at min & 0.0033 & 0.0030 & $0.0000$ & 0.0792 & yes & 0.704 & $+0.2799$ \\
 & rates $\times10$, at min & 0.0109 & 0.0030 & $0.0000$ & 0.1179 & yes & 0.863 & $+0.1149$ \\
 & rates $\times100$, at min & 0.1090 & 0.0030 & $0.0000$ & 0.3346 & yes & 0.962 & $+0.0128$ \\
 & rates $\times3$, $v_t=\sqrt{Q_t}$ & 0.0033 & 0.0030 & $-0.0030$ & 0.0572 & \textbf{no} & 0.975 & $-0.0009$ \\
\midrule
U.S. vs. Germany
 & means, $v_t=\bar v$ & 0.0009 & 0.0030 & $0.0076$ & 0.1070 & yes & 0.263 & $+0.6390$ \\
 & means, $v_t$ min & 0.0009 & 0.0030 & $0.0000$ & 0.0622 & yes & 0.453 & $+0.4762$ \\
 & rates $\times2$, at min & 0.0017 & 0.0030 & $0.0000$ & 0.0689 & yes & 0.578 & $+0.3596$ \\
 & rates $\times3$, at min & 0.0026 & 0.0030 & $0.0000$ & 0.0749 & yes & 0.651 & $+0.2907$ \\
 & rates $\times10$, at min & 0.0087 & 0.0030 & $0.0000$ & 0.1083 & yes & 0.823 & $+0.1252$ \\
 & rates $\times100$, at min & 0.0872 & 0.0030 & $0.0000$ & 0.3004 & yes & 0.938 & $+0.0130$ \\
 & rates $\times3$, $v_t=\sqrt{Q_t}$ & 0.0026 & 0.0030 & $-0.0030$ & 0.0512 & \textbf{no} & 0.954 & $-0.0027$ \\
\bottomrule
\end{tabular}
\end{table}

\section{Minimum incompleteness is an identifying selection rule}\label{app:mincomp}

The exchange-rate dynamics do not separately identify one risk-premium component and the degree of market incompleteness. The paper resolves this by choosing the unidentified quantity, described as chosen ``in a least-squares sense'', to minimise $\sum_t\epsilon_t^2$ subject to $\epsilon_t^2\ge0$ at every date. Under constant pure-currency risk prices the free quantity is the product $\sigma_{z^*}\epsilon_{c^*}$; under time-varying prices it is the correlation $\rho_{zz^*}$.

Although the criterion is expressed as a sum of squared incompleteness shocks, it is affine in the selected reduced-form parameter, because $\epsilon_t^2$ is itself the residual rather than a squared residual in that parameter. The optimisation is therefore a linear programme in one variable, and its solution is a vertex of the feasible set rather than an interior least-squares projection.

The feasible set matters, and two different sets must be distinguished. Write the free quantity as a scalar $c$ entering linearly, $\epsilon_t^2(c)=A_t-\gamma_tc$ for known loadings $\gamma_t$. Then
\begin{align}
\mathcal C_{\text{published}}&=\left\{c:\ \epsilon_t^2(c)\ge0\ \text{for all }t,\ \ c\in\mathcal R\right\},
\label{eq:feasible_published}\\
\mathcal C_{\text{mathematical}}&=\mathcal C_{\text{published}}\cap\left\{c:\ R(c)\succeq0\right\},
\label{eq:feasible_mathematical}
\end{align}
where $\mathcal R$ collects any range restriction on $c$ itself, such as $c\in[-1,1]$ when $c$ is a correlation, and $R(c)\succeq0$ is positive semidefiniteness of the full Brownian correlation matrix of \Cref{app:psd}. The first set collects the constraints the published estimator is described as imposing; the second adds what a valid stochastic model requires. We do not know which the implementation enforced, and \Cref{app:psd_implementation} records that as unknown, so no conclusion below assumes that the reported optimiser respected $\mathcal C_{\text{mathematical}}$.

The proposition is stated for an abstract compact $\mathcal C$, which covers both. Restricting attention to the nonnegativity constraints alone, as an earlier version of this analysis did, can identify the wrong active constraint.

\begin{proposition}[Corner solution of the minimum-incompleteness rule]\label{prop:min_incomplete}
Suppose only that $\mathcal C$ is nonempty and compact and that $\sum_t\gamma_t\ne0$. No convexity of $\mathcal C$ is assumed. Then the problem
\begin{equation}\label{eq:min_incomplete_problem}
\min_{c\in\mathcal C}\ \sum_{t=1}^N\epsilon_t^2(c)
\end{equation}
attains its minimum at a boundary point of $\mathcal C$, at which at least one defining constraint is active; if $\mathcal C$ is a finite union of compact intervals, the minimiser is an endpoint of one component. Moreover:
\begin{enumerate}[label=(\roman*)]
\item If the active constraint at the optimum is one of the nonnegativity constraints, then $\min_t\epsilon_t^2=0$ exactly.
\item If instead the active constraint is a range restriction or the positive-semidefinite boundary, then $\min_t\epsilon_t^2>0$ is possible and the incompleteness constraints need not bind at all.
\item In the constant-risk-price case $\gamma_t\equiv1$ and $c=\sigma_z\epsilon_c+\sigma_{z^*}\epsilon_{c^*}$, the objective is $\sum_tA_t-Nc$, strictly decreasing in $c$. If the largest feasible $c$ is determined by the nonnegativity constraints rather than by $\mathcal R$ or the semidefiniteness requirement, then
\begin{equation}\label{eq:min_incomplete_solution}
c^*=\min_{1\le t\le N}A_t,
\qquad
\epsilon_t^2=A_t-\min_sA_s .
\end{equation}
\end{enumerate}
\end{proposition}

\begin{proof}
The objective equals $\sum_tA_t-c\sum_t\gamma_t$, which is affine in $c$ with slope $-\sum_t\gamma_t\ne0$. An affine function with nonzero slope is strictly monotone, so on a nonempty compact set $\mathcal C\subset\R$ it attains its minimum at $\max\mathcal C$ or at $\min\mathcal C$ according to the sign of the slope; both exist because $\mathcal C$ is compact. Any such point is a boundary point of $\mathcal C$, and at a boundary point at least one defining constraint holds with equality, since if every constraint were slack a neighbourhood would remain feasible. If $\mathcal C$ is a finite union of compact intervals, the minimiser is an endpoint of one component, by the same argument applied to each component. Statements (i) and (ii) follow according to which constraint is active there. For (iii), $\gamma_t\equiv1$ gives slope $-N<0$, so the optimum is $\max\mathcal C$; if that equals $\min_tA_t$ then \eqref{eq:min_incomplete_solution} follows and the constraint indexed by $\arg\min_tA_t$ is active.
\end{proof}

The proposition is conditional in a way the earlier version was not. It does not assert that an incompleteness constraint must always bind, because the optimum may instead be stopped by $\mathcal R$ or by semidefiniteness. Whether the conditional case applies to the published estimates is an empirical question about interiority, and it can be checked only partially. In the time-varying specifications, where Table 4 of \citet{BrandtSantaClara2002} reports $\rho_{zz^*}$ directly, the estimates $0.89$ and $0.94$ are interior to $[-1,1]$; and \Cref{app:psd} finds the reconstructed correlation matrices positive definite with the minimum eigenvalue bounded away from zero at every model-feasible reference state examined. Under those conditions the range and semidefiniteness constraints are slack, so the active constraint is a nonnegativity constraint and case (i) applies. We emphasise that this is an inference from the reported estimates and the states we can reconstruct, not from the observed state series, which we do not have.

Table 4 of \citet{BrandtSantaClara2002} reports the minimum of the inferred excess volatility as $0.0000$ in all six specifications, that is for models A, B and C in both the U.S.--U.K. and the U.S.--Germany systems. This is consistent with case (i), but it should be read carefully: a reported $0.0000$ is a rounded value at four decimal places and is not by itself proof of exact binding. What makes exact binding the right reading here is the optimisation geometry, not the table: the objective is affine with nonzero slope, so the optimum is a vertex, and the interiority checks above locate that vertex on a nonnegativity face. The table is corroboration of a conclusion that follows from \Cref{prop:min_incomplete}, not the evidence for it.

The mechanism is acknowledged in part in the original paper, which observes on p.~193 that with constant risk prices the inferred excess volatility is essentially $\epsilon_t=\{v_t^2-\min_sv_s^2\}^{1/2}$, and draws from this the conclusion that a volatile $v_t$ leaves a large average residual. \Cref{prop:min_incomplete} formalises that observation, shows that it is a property of the optimisation rather than of the data, and extends it to the time-varying specifications where the same corner structure holds for a different free parameter. The consequence for interpretation is that the dispersion of $v_t^2$, rather than its level, mechanically drives much of the estimated excess volatility: the residual is the observed variance minus its sample minimum after other adjustments, so its average size is a measure of that dispersion and is bounded below by construction.

The rule is transparent, but it is not identification by the likelihood of the original structural model. It selects a member of an observationally equivalent or weakly identified class by imposing the smallest feasible incompleteness. The paper describes it as a ``prior of sorts''. Mathematically, the final estimator combines simulated likelihood, auxiliary least squares from bond yields, a minimum-incompleteness corner rule, and inequality constraints. A further consequence, used in \Cref{app:reality}, is that under case (i) the reconstructed incompleteness state sits exactly on its own degenerate boundary at the minimising date.

\section{Binding constraints imply nonregular inference}\label{app:constrained}

The empirical optimisation assigns zero likelihood to parameter values for which any reconstructed $\epsilon_t^2$ is negative. The paper later states that the constraint is binding in the time-varying specifications. A criterion with an active inequality boundary is not twice continuously differentiable across the boundary, and the usual unconstrained normal approximation need not hold. Boundary MLEs and likelihood-ratio statistics commonly have tangent-cone or mixture limits rather than the standard interior distributions; see \citet{SelfLiang1987}.

Stacking first-order conditions and reporting ordinary GMM standard errors does not automatically account for:
\begin{itemize}
\item a sample-dependent feasible set defined by $N$ inequalities;
\item changes in the active set;
\item nondifferentiability of an inner constrained optimisation;
\item structural parameters selected through minimum incompleteness;
\item an optimum on or near the boundary.
\end{itemize}
A separate constrained-estimation analysis is needed before conventional symmetric standard errors and Wald tests can be justified.

\section{The argmax question in a location model}\label{app:argmax}

This appendix records what can be proved about the simulated-likelihood maximiser in the most favourable possible model, and where the argument stops. The conclusion is negative: we do not obtain a counterexample to argmax consistency, and the manuscript makes no such claim. Two intermediate results are complete and are stated here because they clarify what the endpoint construction does to the estimation problem.

Consider
\begin{equation}\label{eq:location_model}
dY_t=\theta\,dt+dW_t,\qquad Y_t\in\R^K,
\end{equation}
observed at unit intervals $t=0,1,\ldots,N$. The Euler scheme is exact for every $M$, the coefficients satisfy Assumptions 1 and 2 of \citet{BrandtSantaClara2002} exactly, and the exact maximum likelihood estimator is the mean increment $\widehat\theta_N=N^{-1}\sum_n\Delta Y_n$. Common random numbers across $\theta$ are used throughout, as the published implementation requires for a smooth objective: the innovations $\xi_{n,s}\sim\Normal(0,I_K)$ are drawn once and held fixed as $\theta$ varies.

\begin{proposition}[Common random numbers give an exact kernel estimator]\label{prop:crn_kde}
Set $h=1/M$, $u_n(\theta)=\Delta Y_n-\theta$ and $a_{n,s}=\sqrt{1-h}\,\xi_{n,s}$. Then for every $\theta$,
\begin{equation}\label{eq:crn_kde}
\widehat q_{M,S}(Y_{n+1}\mid Y_n;\theta)
=\frac1S\sum_{s=1}^S\phik\!\left(u_n(\theta)-a_{n,s};0,hI_K\right).
\end{equation}
That is, the simulated transition density is exactly a Gaussian kernel density estimate with bandwidth $\sqrt h$, built from a sample that does not depend on $\theta$ and evaluated at the $\theta$-dependent point $u_n(\theta)$.
\end{proposition}

\begin{proof}
The Euler recursion for \eqref{eq:location_model} is $Y^M_{m+1}=Y^M_m+\theta h+\sqrt h\,\varepsilon_{m+1}$, so after $M-1$ steps from $Y^M_0=Y_n$ the preterminal state is $Z_{n,s}(\theta)=Y_n+\theta(1-h)+\sqrt{1-h}\,\xi_{n,s}$, where $\sqrt{1-h}\,\xi_{n,s}$ collects the $M-1$ independent increments. The one-step density \eqref{eq:one_step_kernel} evaluated at $Y_{n+1}$ has mean $Z_{n,s}(\theta)+\theta h$ and covariance $hI_K$, and
\[
Y_{n+1}-Z_{n,s}(\theta)-\theta h=\Delta Y_n-\theta-\sqrt{1-h}\,\xi_{n,s}=u_n(\theta)-a_{n,s}.
\]
Substituting into \eqref{eq:pedersen_estimator} gives \eqref{eq:crn_kde}.
\end{proof}

\Cref{prop:crn_kde} is an identity, not an approximation, and it recasts \Cref{thm:collapse} in familiar terms: a kernel density estimate with $S$ atoms and bandwidth $\sqrt h$ is uninformative when $Sh^{K/2}\to0$.

\begin{proposition}[The maximiser is governed by a nearest-atom objective]\label{prop:nearest_atom}
Write $b_{n,s}=\Delta Y_n-\sqrt{1-h}\,\xi_{n,s}$ and
\[
D_N(\theta)=\sum_{n=0}^{N-1}\min_{s\le S}\lVert\theta-b_{n,s}\rVert^2 .
\]
Then for every $\theta$,
\begin{equation}\label{eq:nearest_atom_sandwich}
-2Nh\log S\le
2h\log\widehat{\mathcal L}_{N,M,S}(\theta)+D_N(\theta)+NhK\log(2\pi h)
\le 0 .
\end{equation}
\end{proposition}

\begin{proof}
Fix $n$ and let $d_n=\min_s\lVert\theta-b_{n,s}\rVert^2$. Factoring the minimising term out of the sum in \eqref{eq:crn_kde},
\[
\frac1S\sum_s e^{-\lVert\theta-b_{n,s}\rVert^2/(2h)}
=e^{-d_n/(2h)}\cdot\frac1S\sum_s e^{-(\lVert\theta-b_{n,s}\rVert^2-d_n)/(2h)} .
\]
Every exponent in the second factor is nonpositive and at least one is zero, so that factor lies in $[1/S,1]$ and its logarithm lies in $[-\log S,0]$. Taking logarithms, adding the constant $-\tfrac K2\log(2\pi h)$ from the Gaussian prefactor, multiplying by $2h$ and summing over $n$ gives \eqref{eq:nearest_atom_sandwich}.
\end{proof}

The bound is uniform in $\theta$ with slack $2Nh\log S$, which along $M=S=n$ is $2N(\log n)/n$ per unit of $N$ and therefore negligible relative to $D_N$. The maximiser of the simulated likelihood is thus governed by a nearest-neighbour objective over the $NS$ atoms $b_{n,s}$, whereas the exact log likelihood in this model is the smooth quadratic $-\tfrac12\sum_n\lVert\theta-\Delta Y_n\rVert^2$ with minimiser $\widehat\theta_N$. The uniform bound alone does not imply that the nearest-atom and exact quadratic objectives share either a finite-sample maximiser or the same local asymptotic behaviour; establishing or refuting that requires the limit analysis discussed below, which we do not have. What the bound does establish is that the object being maximised is a nearest-neighbour functional of $NS$ atoms rather than a smooth quadratic, so any argument that the two maximisers agree must be made rather than assumed.

What is missing is the limit of $D_N$. A nearest-neighbour heuristic, given in the accompanying notes, suggests that $N^{-1}S^{2/K}D_N(\theta)$ converges to a multiple of $\exp\{\lVert\theta-\thetao\rVert^2/(K-2)\}$, finite exactly when $K>2$, which is the same threshold as \Cref{thm:collapse}. If that were established, the simulated argmax would be \emph{consistent} for $\thetao$ but would converge at a rate governed by $S^{2/K}$ rather than the parametric rate, so that $\sqrt N(\widehat\theta_{N,M,S}-\widehat\theta_N)$ would not vanish, contradicting the conclusion of Lemma 5. Proving it requires a uniform-in-$\theta$ nearest-neighbour law for a triangular array of exchangeable, not independent, atoms with $S$ growing; standard nearest-neighbour asymptotics are pointwise and assume independent sampling from a fixed density. We did not close that gap.

The status of the ingredients above should be kept distinct.

\Cref{prop:crn_kde} and \Cref{prop:nearest_atom} are an exact identity and a proved bound respectively.

The nearest-neighbour limit is a \emph{conjecture}. It is stated as such in the accompanying notes, it is used nowhere in this paper, and we do not assert that $N^{-1}S^{2/K}D_N(\theta)$ converges to the shape given there. A proof would require uniform-in-$\theta$ nearest-neighbour asymptotics for a triangular array whose atoms are exchangeable rather than independent, since they share the observation innovation, together with control of the dependence that common random numbers induce across $\theta$ and an argmax-continuity argument. None of these is available to us.

The accompanying notes also report \emph{numerical observations}. Across the designs examined the simulated maximiser appears to approach $\thetao$, while the scaled gap $\sqrt N\lVert\widehat\theta_{N,M,S}-\widehat\theta_N\rVert$ shows no downward trend. Those runs use a single seed per design, cover one dimension, and cannot reach the regime $N^4\ll S\ll M^2$ in which the published theorem operates. They are not a demonstration of consistency, and they are deliberately confined to the notes rather than reported as a finding here.

Accordingly the manuscript claims only what \Cref{sec:proof_review} establishes: Lemmas 2 and 3 are directly refuted, and the proofs of Theorems 1 and 2 therefore fail.

\section{Reproducibility}\label{app:repro}

The accompanying directory contains:
\begin{itemize}
\item \texttt{main.tex}, the manuscript source;
\item \texttt{references.bib}, the bibliography;
\item \texttt{code/generate\_results.py}, the fully reproducible simulation and figure code;
\item vector PDF and PNG versions of all figures;
\item CSV and LaTeX versions of generated numerical tables.
\end{itemize}
The code uses a fixed seed, vectorised Gaussian simulation, and log-sum-exp evaluation of the endpoint estimator. No proprietary data are used.

\printbibliography

\end{document}
